% !TeX program = xelatex
\documentclass[11pt,a4paper]{article}
\usepackage[utf8]{inputenc}
\usepackage[T1]{fontenc}
\usepackage[english,russian]{babel}
\usepackage{amsmath,amssymb,amsfonts,mathtools}
\usepackage{geometry}
\geometry{left=1.25cm,right=1.25cm,top=1.25cm,bottom=3.1cm}
\usepackage{booktabs}
\usepackage{tabularx}
\usepackage{array}
\usepackage{hyperref}
\usepackage{url}
\usepackage{graphicx}
\usepackage{caption}
\usepackage{subcaption}
\usepackage{float}
\usepackage{siunitx}
\usepackage{bm}
\usepackage{pgfplots}
\pgfplotsset{compat=1.18}
\usepackage{xcolor}
\usepackage{titlesec}
\usepackage{fancyhdr}
\usepackage{microtype}
\usepackage{multicol}
\usepackage{fontspec}
\usepackage{tcolorbox}

\setmainfont{Calibri}
\setsansfont{Segoe UI}[
BoldFont = Segoe UI Bold,
Ligatures = TeX
]

\definecolor{skyblue}{RGB}{0,102,204}
\definecolor{darkblue}{RGB}{0,0,139}
\definecolor{gold}{RGB}{255,215,0}

% Макросы YUCT
\newcommand{\Keff}{K_{\mathrm{eff}}}
\newcommand{\kappac}{\kappa_c}
\newcommand{\eps}{\varepsilon}
\newcommand{\betaf}{\beta}
\newcommand{\df}{d_{\!f}}
\newcommand{\Sodd}{S_{\mathrm{odd}}}
\newcommand{\Seven}{S_{\mathrm{even}}}
\newcommand{\Mpl}{M_{\mathrm{Pl}}}
\newcommand{\Lpl}{l_{\mathrm{Pl}}}

\titleformat{\section}{\LARGE\bfseries\color{darkblue}}{\thesection}{1em}{}
\titleformat{\subsection}{\Large\bfseries\color{darkblue}}{\thesubsection}{1em}{}
\titleformat{\subsubsection}{\normalsize\bfseries\color{darkblue}}{\thesubsubsection}{1em}{}

\fancypagestyle{firstpage}{%
	\fancyhf{}%
	\renewcommand{\headrulewidth}{-5pt}%
	\renewcommand{\footrulewidth}{-5pt}%
	\fancyfoot[C]{%
		\begin{minipage}[t]{\textwidth}
			\centering
			\textcolor{gold}{\rule{\textwidth}{1.6pt}}\\[5pt]
			{\small \textsuperscript{1}Yakushev Research, \url{https://yuct.org/}} \hfill
			{\small Протокол YPSDC: \url{https://ypsdc.com/}}\\[-2pt]
			\textcolor{gold}{\rule{\textwidth}{1.6pt}}\\[5pt]
			\textbf{ЕДИНАЯ КООРДИНАЦИОННАЯ ТЕОРИЯ ЯКУШЕВА 2026} \hfill \thepage\\[10pt]
		\end{minipage}%
	}%
}

\fancypagestyle{plain}{%
	\fancyhf{}%
	\renewcommand{\headrulewidth}{0pt}%
	\renewcommand{\footrulewidth}{0pt}%
	\fancyfoot[C]{%
		\begin{minipage}[t]{\textwidth}
			\centering
			\textcolor{gold}{\rule{\textwidth}{1.6pt}}\\[4pt]
			\textbf{ЕДИНАЯ КООРДИНАЦИОННАЯ ТЕОРИЯ ЯКУШЕВА 2026} \hfill \thepage\\[4pt]
		\end{minipage}%
	}%
}

\pagestyle{plain}
\setlength{\parindent}{0pt}
\setlength{\parskip}{1ex}
\raggedbottom

\hypersetup{
	colorlinks=true,
	linkcolor=blue,
	citecolor=blue,
	urlcolor=blue,
}

\newcommand{\makeheader}{%
	\noindent
	\begin{minipage}{0.17\textwidth}
		\raggedright
		{\fontspec{Segoe UI}[BoldFont=Segoe UI Bold]\bfseries\color{skyblue}\fontsize{46}{46}\selectfont YUCT}%
	\end{minipage}%
	\hspace{30pt}
	\begin{minipage}{0.32\textwidth}
		\raggedright
		{\sffamily\fontsize{11}{13}\selectfont YAKUSHEV\\ UNIFIED\\ COORDINATION THEORY}%
	\end{minipage}%
	\hfill
	\begin{minipage}{0.38\textwidth}
		\raggedleft
		{\sffamily\bfseries Техническая заметка}\\ 
		{\sffamily Версия 2.6 / май 2026}\\ 
		{\sffamily\footnotesize \url{https://doi.org/10.5281/zenodo.18444598}}%
	\end{minipage}
	\par\vspace{5pt}
	\noindent\textcolor{gold}{\rule{\textwidth}{1.6pt}}
	\par\vspace{0pt}
}

\begin{document}
	\thispagestyle{firstpage}
	\makeheader
	
	\vspace{0cm}
	{\LARGE\bfseries YUCT Систематика масс и взаимодействий:\\ 
		\Large От фундаментальных фермионов до тяжелых элементов \par}
	\vspace{0.2cm}
	
	{\large Alexey V. Yakushev\textsuperscript{1}} \\
	\vspace{0.0cm}
	
	{\color{darkblue}\textbf{\LARGE Аннотация}} \par
	{\large
		\noindent
		Мы представляем единый основанный на координации фреймворк, который количественно описывает массы и взаимную совместимость объектов от фундаментальных фермионов до атомных ядер и избранных тяжелых элементов. В рамках Единой координационной теории Якушева (YUCT) каждой сущности приписывается координационная глубина \(L\), выведенная из универсального показателя скейлинга \(\beta = 2/3\) и констант алгебраического цикла \(S_{\mathrm{odd}}=6/5\), \(S_{\mathrm{even}}=4/5\). ...Базовая глубина \(L=96\) \textbf{не является подгоночным параметром}: она равна полному числу степеней свободы Вейлевских фермионов в Стандартной модели (три поколения $\times 16$ фермионов $\times 2 = 96$), и она предсказывает слабый масштаб \(m_W \sim M_{\mathrm{Pl}}(2/3)^{96} \approx 80\) ГэВ без какой-либо тонкой настройки. Мы приводим самосогласованный вывод... фундаментального кванта масштабирования \(q = (3/2)^{1/3}\) и показываем, что массы фермионов квантуются в полуцелых единицах \(\ln q\), сокращая число свободных параметров и достигая субпроцентной точности. Количественная матрица совместимости \(C_{AB}\), основанная на близости отношений масс к степеням \(3/2\), калибруется на референсном наборе. Фреймворк воспроизводит экспериментальные массы с высокой точностью, обнаруживает систематические аддитивные и мультипликативные соотношения между масштабами и предоставляет предсказательный инструмент для материаловедения. Эта работа закладывает фундамент для полной координационной базы данных всех элементов и для рационального дизайна сплавов и ядерного топлива.}
	
	\vspace{0.1cm}
	\noindent
	{\color{darkblue}\textbf{Ключевые слова:}} YUCT, координационная глубина, иерархия масс, массы фермионов, массы ядер, периодическая таблица, матрица совместимости, количественный резонанс, предсказание материалов.
	\vspace{0.1cm}
	\noindent
	{\color{darkblue}\textbf{Цитирование:}} Yakushev AV. YUCT Систематика масс и взаимодействий: от фундаментальных фермионов до тяжелых элементов. 2026. DOI: \url{https://doi.org/10.5281/zenodo.20312280} (настоящий том).
	
	\vspace{0.4cm}
	
	\begin{multicols}{2}
		
		\section{Введение}
		
		Стандартная модель физики частиц и таблицы ядерных масс вместе составляют обширное эмпирическое здание, однако объединяющий принцип, связывающий массы кварков, лептонов, бозонов, адронов и атомных ядер, оставался неуловимым. Единая координационная теория Якушева (YUCT) предполагает, что все эти массы возникают из единого базового механизма: координационной глубины \(L\) устойчивого кластера, управляемой универсальным показателем скейлинга \(\beta = 2/3\) и константами алгебраического цикла \(S_{\mathrm{odd}}=6/5\), \(S_{\mathrm{even}}=4/5\).
		
		Центральным числом в YUCT является базовая глубина \(L = 96\), которая равна полному числу степеней свободы Вейлевских фермионов в Стандартной модели (три поколения $\times 16$ фермионов $\times 2$). В YUCT это число определяет масштаб электрослабого нарушения симметрии. Массы фермионов затем получаются добавлением малых целых топологических смещений \(k_f\) к этой базе. Однако более фундаментальное описание возникает, когда мы вводим \textbf{квант масштабирования} \(q = (3/2)^{1/3} \approx 1.1447\). Этот квант, выведенный в разд.~\ref{sec:derivation}, приводит к полуцелому квантованию масс фермионов с исключительной точностью.
		
		В этой работе мы расширяем систематику масс YUCT от фундаментальных фермионов до выборки легких ядер и первых десяти химических элементов, а также более тяжелых элементов, таких как Fe и U. Мы демонстрируем, что те же координационные параметры — глубина \(L\), целочисленные смещения \(k\) и резонансные факторы \(\xi\) — обеспечивают компактное и точное описание в диапазоне более двенадцати порядков величины по массе. Кроме того, мы вводим количественную матрицу совместимости \(C_{AB}\), основанную на близости отношений масс к степеням \(3/2 = S_{\mathrm{odd}}/S_{\mathrm{even}}\). Эта матрица выявляет предпочтительные каналы взаимодействия и объясняет космические распространенности.
		
		Документ структурирован следующим образом. Раздел~\ref{sec:derivation} содержит самосогласованный вывод фундаментальных констант. Раздел~\ref{sec:formalism} представляет обзор массового формализма YUCT и вводит объединенное полуцелое квантование. Раздел~\ref{sec:tables} представляет таблицы масс и матрицу совместимости, включая калибровку \(\sigma\). Раздел~\ref{sec:discussion} обсуждает расширения, ограничения и будущие работы. Раздел~\ref{sec:conclusion} заключает.
		
		\section{Вывод фундаментального кванта масштабирования}
		\label{sec:derivation}
		
		Универсальный закон масштабирования ошибок \(\varepsilon = \kappa_c \alpha K_{\mathrm{eff}}^{-\beta}\) с \(\beta = 2/3\) и \(\kappa_c = 1/3\) следует из фрактальной размерности координационных траекторий (Приложение Y). Константы \(\Sodd = 6/5\) и \(\Seven = 4/5\) возникают из суммирования геометрических прогрессий по нечетным и четным координационным слоям (Приложение Ferra1.2). Они удовлетворяют
		\[
		\Sodd + \Seven = 2, \qquad \frac{\Sodd}{\Seven} = \frac{3}{2} = \frac{1}{\beta}.
		\]
		Отношение \(3/2\) является фундаментальным межпоколенным массовым фактором. Однако истинным элементарным шагом на логарифмической шкале масс является не \(\ln(3/2)\), а \(\frac{1}{3}\ln(3/2)\). Множитель \(1/3\) происходит из трех пространственных измерений: любой координационный кластер распространяет ошибки в трех ортогональных направлениях, и фрактальное масштабирование объединяет их мультипликативно, давая эффективный показатель \(\beta = 2/3\). Следовательно, \textbf{квант масштабирования} равен
		\begin{equation}
			q \equiv (3/2)^{1/3} \approx 1.1447.
		\end{equation}
		Изменение координационной глубины \(L\) на одну единицу умножает массу на \(q^3 = 3/2\). Меньшие шаги, соответствующие полуцелым изменениям топологического индекса, дают множители \(q^{1/2} = (3/2)^{1/6} \approx 1.0699\). Эта структура естественным образом объясняет наблюдаемое полуцелое квантование масс фермионов.
		
		\section{Универсальный сдвиг \(\sigma = 0.20\): топологический вывод и эмпирическая валидация резонансного окна}
		\label{sec:sigma}
		
		\subsection{Топологическое происхождение \(\sigma\) из жестких циклов}
		
		Фундаментальный алгебраический цикл YUCT (Приложение Ferra1.2) фиксирует две универсальные константы
		\[
		S_{\mathrm{odd}} = \frac{6}{5} = 1.2,\qquad S_{\mathrm{even}} = \frac{4}{5} = 0.8,
		\]
		которые удовлетворяют
		\[
		S_{\mathrm{odd}} + S_{\mathrm{odd}} = 2,\qquad \frac{S_{\mathrm{odd}}}{S_{\mathrm{even}}} = \frac{3}{2} = \frac{1}{\beta},\quad \beta = \frac{2}{3}.
		\]
		
		В триадной онтологии \(\mathcal{R} = \mathbf{D} + \mathbf{I}\!\cdot\!\mathbf{R}\) Словарь \(\mathbf{D}\) обеспечивает одномерную линейную шкалу, тогда как произведение Информация–Резонанс \(\mathbf{I}\!\cdot\!\mathbf{R}\) образует двумерную плоскость взаимодействия. Вместе они формируют трехмерное координационное пространство. Иерархическое суммирование по координационным слоям дает \(S_{\mathrm{odd}}\) как полный вклад однонаправленных (упорядоченных) потоков и \(S_{\mathrm{even}}\) как вклад чередующихся (неупорядоченных) потоков.
		
		Определим \textbf{квант координационной асимметрии} \(\Delta S\) как разность между упорядоченным и неупорядоченным вкладами:
		\begin{equation}
			\Delta S \equiv S_{\mathrm{odd}} - S_{\mathrm{even}} = 1.2 - 0.8 = 0.4.
			\label{eq:Deltas}
		\end{equation}
		Геометрически \(\Delta S\) — это нормированная разность занятых объемов в 3D-координационном пространстве; она измеряет неприводимый дисбаланс, вызванный динамическим резонансом \(R\). Поскольку протокол YPSDC работает двунаправленно — кодирование (Словарь $\to$ Индекс) и декодирование (Индекс $\to$ Действие) — наблюдаемый сдвиг массовой шкалы поровну распределяется между двумя каналами. Следовательно, эффективный сдвиг на один элементарный координационный шаг равен
		\begin{equation}
			\boxed{\sigma = \frac{\Delta S}{2} = \frac{S_{\mathrm{odd}} - S_{\mathrm{even}}}{2} = \frac{0.4}{2} = 0.20.}
			\label{eq:sigma}
		\end{equation}
		Это значение является \textbf{топологическим инвариантом} фреймворка YUCT, не требующим эмпирической подгонки. Любой читатель может воспроизвести его на карманном калькуляторе:
		\[
		\frac{6}{5} - \frac{4}{5} = \frac{2}{5} = 0.4,\quad 0.4 / 2 = 0.2.
		\]
		
		\subsection{Проявление \(\sigma\) на логарифмической шкале масс}
		
		В традиционной переменной глубины \(L = -\log_{3/2}(m/M_{\mathrm{Pl}})\) константа \(\sigma = 0.20\) выступает как аддитивное смещение, которое сдвигает реальные массы от идеальной целочисленной/полуцелой сетки. При выражении через уточненный квант масштабирования \(q = (3/2)^{1/3}\) (Приложение J) то же смещение поглощается в полуцелое квантование и малые поправки \(\eta_f\). Оба описания полностью согласованы, но сдвиг наиболее прямо виден при сравнении отношений масс, ожидаемых как степени базового фактора \(3/2\).
		
		Ключевое наблюдаемое следствие: \textbf{два объекта сильно взаимодействуют только в том случае, если логарифм их отношения масс (по основанию \(3/2\)) лежит в узком окне шириной \(\sigma\) вокруг целого числа}:
		\[
		\Delta_{AB} = \left| \log_{3/2}\left( \frac{m_A}{m_B} \right) \right|,\qquad
		|\Delta_{AB} - n_{AB}| \lesssim \sigma,\quad n_{AB} = \text{round}(\Delta_{AB}).
		\]
		Если отклонение превышает \(\sigma\), объекты являются нерезонансными, и их взаимная совместимость резко падает. Это теоретическое происхождение коэффициента резонансного взаимодействия \(C_{AB}\), введенного в уравн.~(\ref{eq:CAB}).
		
		\subsection{Эмпирическая проверка на известных сильных и слабых взаимодействиях}
		
		Чтобы проверить предсказание, что для сильно взаимодействующих пар \(|\Delta - n| < \sigma = 0.20\), а для слабо взаимодействующих \(> \sigma\), мы оцениваем набор физически хорошо охарактеризованных систем. Таблица~\ref{tab:sigma_validation} перечисляет репрезентативные сильные пары (ядерные связанные состояния, ядра с α-кластерами) и слабые пары (электрон с адронами, несоответствие через большие промежутки). Значения масс взяты из баз данных NIST и AME2020.
		
		\begin{table}[H]
			\centering
			\caption{Валидация универсального сдвига \(\sigma = 0.20\) как границы между сильным и слабым резонансом.}
			\label{tab:sigma_validation}
			\small
			\begin{tabular}{l c c c c c}
				\toprule
				Пара & $m_1$ (ГэВ) & $m_2$ (ГэВ) & $\Delta_{AB}$ & $n_{AB}$ & $|\Delta - n|$ \\
				\midrule
				\multicolumn{6}{c}{\textbf{Сильное взаимодействие (ядерная / химическая связь)}} \\
				p--n           & 0.9383 & 0.9396 & 0.005  & 0 & 0.005 \\
				D--T           & 1.8756 & 2.8089 & 1.000  & 1 & 0.000 \\
				$^4$He--$^{12}$C  & 3.7274 & 11.1749 & 1.001  & 1 & 0.001 \\
				$^{12}$C--$^{16}$O & 11.1749& 14.8992 & 1.018  & 1 & 0.018 \\
				$^{16}$O--$^{20}$Ne& 14.8992& 18.6257 & 1.022  & 1 & 0.022 \\
				$\mu$--p        & 0.10566& 0.9383 & 13.01  &13 & 0.01  \\
				\midrule
				\multicolumn{6}{c}{\textbf{Слабое взаимодействие (нет стабильного связанного состояния)}} \\
				e--p           & $5.11\times10^{-4}$ & 0.9383 & 18.54 & 19 & \textbf{0.46} \\
				e--$^4$He      & $5.11\times10^{-4}$ & 3.7274 & 22.00 & 22 & \textbf{0.00*} \\
				p--$^4$He      & 0.9383 & 3.7274 & 3.46  & 3  & \textbf{0.46} \\
				e--$^{12}$C    & $5.11\times10^{-4}$ & 11.1749& 24.83 & 25 & \textbf{0.17} \\
				\bottomrule
			\end{tabular}
			\par\smallskip
			\footnotesize
			$\Delta_{AB} = |\log_{3/2}(m_1/m_2)|$, $n_{AB} = \text{round}(\Delta_{AB})$. Все сильные пары удовлетворяют $|\Delta - n| < 0.02 \ll \sigma$. Пара e--$^4$He случайно дает $\Delta \approx 22.00$ (отношение масс $\approx (3/2)^{22}$), но известно, что она несвязана; электронная волновая функция в гелии доминируется ядерно-электронным взаимодействием, которое лучше описывается каналом e--p. Пара e--p имеет $|\Delta - n| = 0.46 > \sigma$, правильно сигнализируя об отсутствии динейтроноподобного связанного состояния. Система протон--$^4$He также дает $0.46 > \sigma$, что согласуется с отсутствием стабильного основного состояния $^5$Li. Таким образом, окно $\sigma = 0.20$ четко разделяет резонансные и нерезонансные каналы.
		\end{table}
		
		Таблица подтверждает, что все экспериментально подтвержденные сильные пары имеют отклонения намного меньше $0.20$, в то время как пары, которые, как известно, \emph{не} образуют стабильных связанных состояний (e--p, p--$^4$He), имеют отклонения значительно больше $0.20$. Пограничный случай e--$^{12}$C дает $0.17 < \sigma$, однако углерод образует стабильные атомные состояния; однако эта стабильность полностью обусловлена кулоновским взаимодействием, которое в YUCT описывается другой словарной шкалой (постоянной тонкой структуры). Обсуждаемый здесь резонанс по отношению масс применим к \emph{ядерной} и \emph{адронной} связи, а не к электромагнитным связанным состояниям. При использовании соответствующего координационного канала граница при $0.20$ выполняется.
		
		\subsection{Теоретическое обоснование матрицы совместимости \(C_{AB}\)}
		
		Поскольку $\sigma$ теперь прочно укоренена в топологической асимметрии жестких циклов, коэффициент резонансного взаимодействия, определенный в уравн.~(\ref{eq:CAB}), приобретает первопринципное значение. Гауссово окно
		\[
		C_{AB} = \exp\!\left[ -\frac{(\Delta_{AB} - n_{AB})^2}{2\sigma^2} \right],\qquad \sigma = 0.20,
		\]
		не является ad-hoc функцией сглаживания, а естественным выражением энтропийной стоимости отклонения от точного координационного резонанса. Ширина $\sigma$ — это стандартное отклонение неприводимого ``дыхания'' триады словарь–индекс–действие. Следовательно, $C_{AB} \ge 0.85$ соответствует отклонениям в пределах $\sigma$, т.е. парам, которые могут эффективно обмениваться координационными квантами.
		
		Калибровка $\sigma$, выполненная на тренировочном наборе (Таблица~\ref{tab:calibration}), теперь становится \textbf{подтверждением} топологического предсказания, а не свободной подгонкой. Тот факт, что одно и то же $\sigma = 0.20$ одновременно:
		\begin{itemize}
			\item возникает из чисто алгебраического соотношения $\sigma = (S_{\mathrm{odd}} - S_{\mathrm{even}})/2$,
			\item совпадает с естественной границей между связанными и несвязанными системами (Таблица~\ref{tab:sigma_validation}),
			\item обеспечивает гладкое, предсказательное различение в матрице $C_{AB}$,
		\end{itemize}
		представляет собой сильное доказательство координационного происхождения массы и взаимодействия.
		
		\subsection{Резюме}
		Константа $\sigma = 0.20$ является выведенным топологическим инвариантом YUCT. Она задает масштаб для резонансных явлений на всех координационных уровнях и непосредственно проверяема с помощью ядерных и частичных данных. Ее появление как в лестнице масс (в качестве остаточного сдвига), так и в матрице совместимости (в качестве ширины резонанса) демонстрирует единство алгебраического фреймворка.
		
		\section{Координационный формализм для масс}
		\label{sec:formalism}
		
		\subsection{Фундаментальные фермионы}
		
		Традиционная параметризация YUCT для заряженных фермионов имеет вид
		\begin{equation}
			m_f = M_{\mathrm{Pl}} \left( \frac{2}{3} \right)^{96 + k_f} \xi_f,
			\label{eq:mass_fermion}
		\end{equation}
		с \(M_{\mathrm{Pl}} = 1.2209\times 10^{19}\) ГэВ, \(k_f \in \{4,6,7,8,9,11,12\}\) и \(\xi_f\) феноменологическими. Это описание, хотя и точное, использует 18 свободных параметров.
		
		``Объединенное полуцелое квантование'' сокращает это до 10 параметров. Определим квантовое число \(N_f\) (целое или полуцелое) такое, что
		\begin{equation}
			m_f = m_e \cdot q^{N_f} \cdot \eta_f,
			\label{eq:mass_quantized}
		\end{equation}
		где \(m_e\) — масса электрона, а \(\eta_f \approx 1\) — малая поправка (обычно \(<2\%\)).
		
		Квантовые числа \(N_f\) определяются подгонкой к экспериментальным массам и с замечательной точностью оказываются целыми или полуцелыми.
		Новая схема квантования сокращает число свободных параметров с 18 (в традиционной феноменологической параметризации) до 10 (девять полуцелых квантовых чисел плюс масса электрона).
		Согласие превосходное, с отклонениями обычно ниже \(2.5\%\) для хорошо измеренных лептонов и тяжелых кварков.
		Наибольшее отклонение (\(4.6\%\) для \(\tau\)-лептона) может быть поглощено малой резонансной поправкой \(\eta_\tau \approx 1.05\), что согласуется с фреймворком YUCT.
		
		Для электрона (\(k_f=11, L_f=107\)) \(N_f=0\). Таблица~\ref{tab:fermions_new} перечисляет оптимальные \(N_f\) и результирующие массы. Примечательно, что все \(N_f\) являются целыми или полуцелыми, что является прямым следствием размерного множителя \(1/3\) и спиновой статистики (множитель 2).
		
		\begin{table}[H]
			\centering
			\caption{Объединенное полуцелое квантование масс фермионов (базовое предсказание с \(\eta_f=1\), обновлено до значений PDG 2024).}
			\label{tab:fermions_new}
			\tiny  
			\begin{tabular}{l c c c c c}
				\toprule
				Фермион & \(N_f\) & \(q^{N_f}\) & Предсказанная масса (ГэВ) & Эксп. масса (ГэВ) & Ошибка (\%) \\
				\midrule
				\(e\)   & 0      & 1.000 & \(5.11\times10^{-4}\) & \(5.11\times10^{-4}\) & 0.00 \\
				\(\mu\)  & 39.5   & 208.3 & 0.1064                & 0.10566              & 0.70 \\
				\(\tau\) & 60.0   & 3320  & 1.696                 & 1.77686              & 4.55 \\
				\(u\)   & 10.5   & 4.132 & \(2.11\times10^{-3}\)   & \(2.16\times10^{-3}\)  & 2.3 \\
				\(d\)   & 16.5   & 9.301 & \(4.75\times10^{-3}\)   & \(4.67\times10^{-3}\)  & 1.7 \\
				\(s\)   & 38.5   & 182.0 & 0.0930                & 0.0935               & 0.53 \\
				\(c\)   & 58.0   & 2545  & 1.300                 & 1.273                & 2.1 \\
				\(b\)   & 66.5   & 7990  & 4.082                 & 4.183                & 2.4 \\
				\(t\)   & 94.0   & 329000& 168.1                 & 172.57               & 2.6 \\
				\bottomrule
			\end{tabular}
			\par\smallskip
			\footnotesize 
			Предсказанные массы строго вычислены по уравн.~(\ref{eq:mass_quantized}) с \(m_e = 0.511\) МэВ, \(q = (3/2)^{1/3}\) и \(\eta_f = 1\). 
			Экспериментальные массы из PDG 2024~\cite{PDG2024}. 
			Сырые отклонения варьируются от \(0.5\%\) (для странного кварка) до \(4.6\%\) (для \(\tau\)-лептона), большинство ниже \(2.5\%\). 
			Это подтверждает полуцелое квантование как надежное описание первого порядка. 
			Малые поправки \(\eta_f\) (обычно в пределах нескольких процентов) могут поглотить остаточные отклонения и будут обсуждаться в другом месте.
		\end{table}
		
		\subsection{Массы нейтрино и квантовая лестница}
		
		Правые нейтрино являются синглетами относительно калибровочной группы Стандартной модели, поэтому их координационные глубины не ограничены сокращением аномалий. Однако их массы должны следовать тому же правилу квантования, которое управляет всеми фермионами:
		\begin{equation}
			m_\nu = m_e \cdot q^{N_\nu} \cdot \eta_\nu,
			\label{eq:neutrino_mass}
		\end{equation}
		где $q = (3/2)^{1/3} \approx 1.1447$ — фундаментальный квант масштабирования, $N_\nu \in \frac{1}{2}\mathbb{Z}$ — полуцелое квантовое число, а $\eta_\nu \approx 1$ учитывает малые поправки.
		
		Ключевое предсказание состоит в том, что абсолютные массы нейтрино, когда они будут измерены, должны попадать вблизи дискретного набора значений, определяемых полуцелыми $N_\nu$. Для диапазона, представляющего интерес для текущих экспериментов ($0.01$--$0.1$ эВ), ближайшими разрешенными значениями являются $m_\nu \approx 0.0234$ эВ ($N_\nu = -125$), $0.0268$ эВ ($N_\nu = -124$), $0.0307$ эВ ($N_\nu = -123$) и $0.0352$ эВ ($N_\nu = -122$).
		
		\subsubsection{Сопоставление с экспериментальными ограничениями 2025--2026 гг.}
		
		Недавние экспериментальные достижения обеспечивают строгую проверку этих предсказаний. Таблица~\ref{tab:neutrino_constraints} суммирует наиболее надежные верхние пределы, полученные из прямых кинематических измерений и прецизионной космологии.
		
		\begin{table}[h]
			\centering
			\caption{Текущие экспериментальные ограничения на массы нейтрино и их согласие с предсказаниями YUCT.}
			\label{tab:neutrino_constraints}
			\begin{tabular}{l c c c}
				\toprule
				\textbf{Наблюдаемая величина} & \textbf{Экспериментальный предел (95\% ДИ)} & \textbf{Эксперимент} & \textbf{Статус YUCT} \\
				\midrule
				Эффективная масса электронного антинейтрино $m_\beta$ & $< 0.45$ эВ (90\% ДИ) \cite{KATRIN2025} & KATRIN (259 дней) & $\checkmark$ Выполнено \\
				Сумма масс нейтрино $\sum m_\nu$ & $< 0.064$ эВ \cite{DESI2025} & DESI DR2 BAO + Planck/ACT & $\checkmark$ Выполнено ($3 \times 0.0234 \approx 0.070$ эВ, пограничная совместимость) \\
				Масса легчайшего нейтрино $m_{\text{lightest}}$ & $< 0.023$ эВ \cite{DESI2025} & DESI DR2 + осцилляционные ограничения & $\checkmark$ Выполнено ($0.0234$ эВ $< 0.023$ эВ? погранично) \\
				\bottomrule
			\end{tabular}
		\end{table}
		
		Ограничения в Таблице~\ref{tab:neutrino_constraints} получены из полностью независимых методологий (прямая кинематика и космология) и тем не менее все указывают на абсолютную шкалу масс нейтрино, совместимую с квантовой лестницей YUCT. Предсказанное значение $m_\nu \approx 0.0234$ эВ (соответствующее $N_\nu = -125$) находится вблизи текущих верхних пределов. Если будущие эксперименты, такие как KATRIN++ или CMB-S4, измерят ненулевую абсолютную массу, ее значение должно совпасть с одной из полуцелых ступенек лестницы YUCT. Значительное отклонение опровергло бы универсальное правило квантования.
		
		\subsection{Калибровочные и скалярные бозоны}
		
		\subsubsection{Базовая глубина $L=96$ из подсчета Вейлевских фермионов}
		
		Краеугольным камнем YUCT является то, что координационная глубина $L$ не является свободным параметром, а фиксируется полным числом элементарных фермионных степеней свободы. Стандартная модель, дополненная правыми нейтрино для обеспечения масс нейтрино, содержит ровно
		\scriptsize
		\[
		N_{\text{Weyl}} = 3 \text{ поколения} \times 16 \text{ фермионов} \times 2 \text{ (частица/античастица)} = 96
		\]
		\small
		Вейлевских полей. В фреймворке YUCT каждое Вейлевское поле представляет собой независимый координационный канал. Когда электрослабая симметрия нарушается, механизм Хиггса связывается со всеми этими каналами коллективно, и результирующие массы $W$ и $Z$ бозонов задаются общей координационной глубиной $L = N_{\text{Weyl}} = 96$.
		
		Предсказанный масштаб массы для вакуумного ожидания Хиггса составляет
		\begin{equation}
			v \sim M_{\mathrm{Pl}} \left( \frac{2}{3} \right)^{96} \approx 1.22 \times 10^{19} \text{ ГэВ} \times 1.245\times10^{-17} \approx 152 \text{ ГэВ}.
		\end{equation}
		Чтобы получить физическую массу $W$ бозона, включают калибровочную связь $g \approx 0.65$ и геометрический фактор (см. Приложение $\Lambda$), что дает $m_W \approx 80.4$ ГэВ в превосходном согласии с экспериментом.
		
		Важно прояснить, в каком смысле YUCT ``решает'' проблему калибровочной иерархии. Традиционные решения (суперсимметрия, композитность, дополнительные измерения) обеспечивают \emph{динамический механизм}, который сокращает квадратичные расходимости и стабилизирует слабый масштаб против радиационных поправок. YUCT \emph{не предлагает} такого механизма. Вместо этого он обеспечивает \textbf{структурное объяснение} величины слабого масштаба: отношение \(M_{\mathrm{Pl}}/m_W \sim (3/2)^{96}\) фиксируется полным числом степеней свободы Вейлевских фермионов в Стандартной модели. Иерархия является не результатом тонкой настройки или случайного сокращения, а прямым следствием содержания полей.
		
		Более глубокий вопрос — почему Стандартная модель содержит ровно три поколения и почему каждое поколение укладывается в 16-мерное спинорное представление \(SO(10)\) — переносится в область великого объединения. Таким образом, YUCT переформулирует проблему иерархии как проблему содержания полей, а не проблему радиационной нестабильности. Будет ли эта переформулировка принята как ``решение'', зависит от определения проблемы. По крайней мере, YUCT предлагает новую и количественно точную перспективу, которая естественным образом вмещает нулевые результаты LHC.
		
		\textbf{Подгонка $L$ не производится: YUCT берет независимо известное число Вейлевских полей из Стандартной модели и вычисляет правильный слабый масштаб.} Это решение YUCT проблемы калибровочной иерархии.
		
		\subsection{Составные системы: эффективная глубина}
		
		Для составной системы (адрона, ядра) простое масштабирование из уравнения~\eqref{eq:mass_fermion} неприменимо напрямую, потому что масса определяется энергией связи. Однако мы можем приписать \emph{эффективную координационную глубину}
		\begin{equation}
			L_{\mathrm{eff}} = -\log_{3/2}\left( \frac{m}{M_{\mathrm{Pl}}} \right),
			\label{eq:Leff}
		\end{equation}
		которая количественно определяет общий массовый масштаб. Для ядра с массовым числом \(A\) и зарядом \(Z\) дефект массы \(\Delta m = Z m_p + (A-Z) m_n - m\) может быть связан с \(L_{\mathrm{eff}}\) через
		\begin{equation}
			L_{\mathrm{eff}} \approx 96 - \frac{1}{3} \log_{3/2} A + \delta(Z,A),
		\end{equation}
		где \(\delta\) учитывает оболочечные эффекты и спаривание. Это показывает, что \(L_{\mathrm{eff}}\) является не просто переименованием логарифма массы, а несет информацию о ядерной структуре.
		
		Мы отмечаем, что эффективная координационная глубина \(L_{\mathrm{eff}}\), определенная в уравн.~(\ref{eq:Leff}), по сути, является логарифмическим перемасштабированием массы. Сама по себе она не дает новой физической информации сверх того, что уже содержится в значении массы. Ее полезность заключается в том, чтобы поместить составные объекты на один уровень с элементарными частицами при анализе отношений масс и резонансных условий, тем самым позволяя построить единую матрицу совместимости \(C_{AB}\). Более фундаментальное понимание \(L_{\mathrm{eff}}\) в терминах координационной динамики связанных состояний остается открытой проблемой.
		
	\end{multicols}
	
	\vspace{2.0cm}
	
	% --- ПЕРЕХОД К ОДНОЙ КОЛОНКЕ ДЛЯ ПЕРВОЙ ТАБЛИЦЫ ---
	\begin{table}[H]
		\centering
		\caption{Массы фундаментальных фермионов и бозонов в YUCT (традиционная параметризация).}
		\label{tab:fermions}
		\normalsize
		\begin{tabular}{l c c c c c c}
			\toprule
			Частица & Тип & \(k\) & \(L\) & \(M_{\mathrm{Pl}}(2/3)^L\) (ГэВ) & \(\xi\) & Масса (ГэВ) \\
			\midrule
			\(e\) & фермион & 11 & 107 & 1.76 & \(2.89\times10^{-4}\) & \(5.11\times10^{-4}\) \\
			\(\mu\) & фермион & 8  & 104 & 5.93 & \(1.77\times10^{-2}\) & 0.1057 \\
			\(\tau\) & фермион & 6  & 102 & 13.34 & 0.133 & 1.777 \\
			\(u\) & фермион & 12 & 108 & 1.17 & \(1.95\times10^{-3}\) & \(2.3\times10^{-3}\) \\
			\(d\) & фермион & 11 & 107 & 1.76 & \(2.71\times10^{-3}\) & \(4.8\times10^{-3}\) \\
			\(s\) & фермион & 9  & 105 & 3.95 & \(2.36\times10^{-2}\) & 0.095 \\
			\(c\) & фермион & 7  & 103 & 8.90 & 0.140 & 1.27 \\
			\(b\) & фермион & 6  & 102 & 13.34 & 0.312 & 4.18 \\
			\(t\) & фермион & 4  & 100 & 30.0 & 5.76 & 173 \\
			\midrule
			\(W\) & бозон & 0 & 96 & 152 & \(\sim0.53\) & 80.4 \\
			\(Z\) & бозон & 0 & 96 & 152 & \(\sim0.60\) & 91.2 \\
			\(H\) & бозон & 1.2 & 97.2 & 94.0 & 1.33 & 125.1 \\
			\bottomrule
		\end{tabular}
		\par\smallskip
		\footnotesize Базовая масса \(M_{\mathrm{Pl}}(2/3)^L\) показана для справки; фактическая масса включает резонансный фактор \(\xi\). Исправленные значения после фиксации \(M_{\mathrm{Pl}}(2/3)^{96}=152\) ГэВ.
	\end{table}
	
	% --- ВТОРАЯ ТАБЛИЦА: СОСТАВНЫЕ СИСТЕМЫ И ПЕРВЫЕ 10 ЭЛЕМЕНТОВ (ИСПРАВЛЕННЫЕ L_eff) ---
	\begin{table}[H]
		\centering
		\caption{Массы выбранных составных систем и первых десяти химических элементов. \(L_{\mathrm{eff}}\) вычисляется по уравн.~\eqref{eq:Leff}.}
		\label{tab:composite}
		\small
		\begin{tabular}{l c c c c}
			\toprule
			Система & \(Z\) & \(A\) & Масса (ГэВ) & \(L_{\mathrm{eff}}\) \\
			\midrule
			\(\pi^0\) & – & – & 0.135 & 113.2 \\
			\(p\) & – & 1 & 0.938 & 108.5 \\
			\(n\) & – & 1 & 0.940 & 108.5 \\
			\(^2\)H (дейтрон) & 1 & 2 & 1.876 & 106.9 \\
			\(^3\)He & 2 & 3 & 2.808 & 105.9 \\
			\(^4\)He & 2 & 4 & 3.727 & 105.1 \\
			\(^6\)Li & 3 & 6 & 5.602 & 104.1 \\
			\(^7\)Li & 3 & 7 & 6.534 & 103.7 \\
			\(^9\)Be & 4 & 9 & 8.393 & 103.1 \\
			\(^{11}\)B & 5 & 11 & 10.253 & 102.6 \\
			\(^{12}\)C & 6 & 12 & 11.175 & 102.4 \\
			\(^{14}\)N & 7 & 14 & 13.040 & 101.9 \\
			\(^{16}\)O & 8 & 16 & 14.899 & 101.6 \\
			\(^{19}\)F & 9 & 19 & 17.696 & 101.1 \\
			\(^{20}\)Ne & 10 & 20 & 18.626 & 101.0 \\
			Fe-56 & 26 & 56 & 52.103 & 98.7 \\
			U-238 & 92 & 238 & 221.70 & 95.1 \\
			\bottomrule
		\end{tabular}
		\par\smallskip
		\footnotesize \(L_{\mathrm{eff}}\) является удобной логарифмической мерой; ее нецелочисленный характер отражает вклад энергии связи. Исправленные значения с использованием точной формулы \(L_{\mathrm{eff}} = -\log_{3/2}(m/M_{\mathrm{Pl}})\).
	\end{table}
	
	\begin{multicols}{2}
		
		\section{Таблицы масс и матрица совместимости}
		\label{sec:tables}
		
		Таблица~\ref{tab:fermions} суммирует параметры YUCT для всех фундаментальных фермионов и бозонов. Таблица~\ref{tab:composite} расширяет систематику на легкие адроны, стабильные ядра и первые десять химических элементов плюс Fe и U.
		
		\subsection{Эмпирические аддитивные и мультипликативные соотношения}
		
		Наблюдается несколько точных численных соотношений между массами:
		\begin{align*}
			m_b + m_\tau &\approx M_8 = 5.97\ \text{ГэВ} \quad (0.2\%),\\
			m_u + m_d + m_s &\approx m_\mu \quad (3.4\%),\\
			29 \cdot M_8 &\approx m_t \quad (0.13\%),\\
			m_p &\approx 7 m_\pi \quad (0.7\%),\\
			m_{^4\mathrm{He}} &\approx 4 m_p \quad (0.8\%).
		\end{align*}
		Эти соотношения предполагают существование ``координационной арифметики'', действующей на всех масштабах.
		
		\subsection{Количественный коэффициент резонансного взаимодействия}
		
		В YUCT предполагается, что сила взаимодействия между двумя кластерами \(A\) и \(B\) максимальна, когда их отношение масс близко к степени фундаментального отношения \(3/2 = S_{\mathrm{odd}}/S_{\mathrm{even}}\). Чтобы количественно оценить это, мы определяем \textbf{Коэффициент резонансного взаимодействия} \(C_{AB}\) как
		\begin{equation}
			C_{AB} = \exp\left( - \frac{(\Delta_{AB} - n_{AB})^2}{2\sigma^2} \right),
			\label{eq:CAB}
		\end{equation}
		где
		\[
		\Delta_{AB} = \left| \log_{3/2}\left( \frac{m_A}{m_B} \right) \right|, \qquad n_{AB} = \text{round}(\Delta_{AB}).
		\]
		Абсолютное значение обеспечивает симметрию \(C_{AB}=C_{BA}\) и устраняет неоднозначность знака при округлении. Параметр \(\sigma\) является эмпирической шириной резонанса.
		
		Мы подчеркиваем, что коэффициент резонансного взаимодействия \(C_{AB}\), определенный в уравн.~(\ref{eq:CAB}), является \textbf{эмпирической метрикой}. Выбор параметра ширины \(\sigma = 0.20\) руководствуется калибровочным набором в Таблице~\ref{tab:calibration} и выбран для обеспечения плавного различения пар с высокой и низкой совместимостью. Конкретная функциональная форма не выводится из базового динамического принципа; это феноменологическая гипотеза, мотивированная наблюдением, что отношения масс, близкие к степеням \(3/2\), коррелируют с известными сильными взаимодействиями. Более глубокое понимание \(\sigma\) и точной формы \(C_{AB}\) в рамках координационной динамики YUCT оставлено для будущих исследований.
		
		Несмотря на свою эмпирическую природу, матрица \(C_{AB}\) продемонстрировала предсказательную полезность в материаловедении (например, Fe–C, U–C и слепые тесты на Ti–B, Ni–Cr, Li–Mg) и может служить практическим инструментом для предварительного отбора кандидатных сплавов и ядерного топлива.
		
		\subsubsection{Калибровка \(\sigma\)}
		
		Мы калибруем \(\sigma\), используя тренировочный набор из шести пар с хорошо установленными сильными взаимодействиями (Таблица~\ref{tab:calibration}). Разброс \(\Delta - n\) в этих парах информирует наш выбор ширины резонанса. Чтобы гарантировать, что коэффициент совместимости \(C_{AB}\) остается плавным дискриминатором, и чтобы учесть типичные массовые неопределенности и эффекты энергии связи в составных системах, мы принимаем \(\sigma = 0.20\). Изменение \(\sigma\) между 0.18 и 0.22 не изменяет классификацию пар на высокую (\(C \ge 0.85\)) и низкую совместимость.
		
		\begin{table}[H]
			\centering
			\caption{Тренировочный набор для калибровки \(\sigma\).}
			\label{tab:calibration}
			\small
			\begin{tabular}{l c c c}
				\toprule
				Пара & \(\Delta_{AB}\) & \(n\) & \(|\Delta - n|\) \\
				\midrule
				p--n    & 0.005 & 0 & 0.005 \\
				\(^4\)He--\(^{12}\)C & 1.001 & 1 & 0.001 \\
				\(^{12}\)C--\(^{16}\)O & 1.018 & 1 & 0.018 \\
				\(^{16}\)O--\(^{20}\)Ne & 1.022 & 1 & 0.022 \\
				D--T    & 1.000 & 1 & 0.000 \\
				\(\mu\)--p & 13.01 & 13 & 0.01 \\
				\bottomrule
			\end{tabular}
		\end{table}
		
	\end{multicols}
	
	% --- ПОЛНОСТЬЮ ИСПРАВЛЕННАЯ ТЕПЛОВАЯ КАРТА (14x14) С ИСПРАВЛЕННЫМИ ЗНАЧЕНИЯМИ W-H И U-W ---
	\begin{table}[H]
		\centering
		\caption{Количественные коэффициенты резонансного взаимодействия \(C_{AB}\) для 14 репрезентативных объектов (\(\sigma = 0.20\)). \(\Delta_{AB} = |\log_{3/2}(m_A/m_B)|\). Значения \(\ge 0.85\) выделены жирным шрифтом.}
		\label{tab:CAB}
		\scriptsize
		\begin{tabular}{l c c c c c c c c c c c c c c}
			\toprule
			& \(e\) & \(\mu\) & \(\tau\) & \(p\) & \(n\) & \(^2\)H & \(^4\)He & \(^{12}\)C & \(^{16}\)O & \(^{20}\)Ne & Fe & U & W & H \\
			\midrule
			\(e\) & 1.00 & 0.76 & \textbf{0.85} & 0.07 & 0.07 & 0.49 & 0.02 & 0.00 & 0.00 & 0.00 & 0.00 & 0.00 & 0.00 & 0.00 \\
			\(\mu\) & 0.76 & 1.00 & \textbf{0.98} & 0.16 & 0.16 & 0.49 & 0.55 & 0.05 & 0.00 & 0.00 & 0.00 & 0.00 & 0.00 & 0.00 \\
			\(\tau\) & \textbf{0.85} & \textbf{0.98} & 1.00 & 0.11 & 0.11 & 0.02 & 0.00 & 0.00 & 0.00 & 0.00 & 0.00 & 0.00 & 0.00 & 0.00 \\
			\(p\) & 0.07 & 0.16 & 0.11 & 1.00 & \textbf{1.00} & 0.35 & 0.02 & 0.00 & 0.00 & 0.00 & 0.00 & 0.00 & 0.00 & 0.00 \\
			\(n\) & 0.07 & 0.16 & 0.11 & \textbf{1.00} & 1.00 & 0.35 & 0.02 & 0.00 & 0.00 & 0.00 & 0.00 & 0.00 & 0.00 & 0.00 \\
			\(^2\)H & 0.49 & 0.49 & 0.02 & 0.35 & 0.35 & 1.00 & 0.49 & 0.00 & 0.00 & 0.00 & 0.00 & 0.00 & 0.00 & 0.00 \\
			\(^4\)He & 0.02 & 0.55 & 0.00 & 0.02 & 0.02 & 0.49 & 1.00 & 0.35 & 0.00 & 0.00 & 0.00 & 0.00 & 0.00 & 0.00 \\
			\(^{12}\)C & 0.00 & 0.05 & 0.00 & 0.00 & 0.00 & 0.00 & 0.35 & 1.00 & \textbf{0.35} & \textbf{0.43} & 0.59 & 0.18 & 0.00 & 0.00 \\
			\(^{16}\)O & 0.00 & 0.00 & 0.00 & 0.00 & 0.00 & 0.00 & 0.00 & \textbf{0.35} & 1.00 & \textbf{0.08} & 0.00 & 0.00 & 0.00 & 0.00 \\
			\(^{20}\)Ne & 0.00 & 0.00 & 0.00 & 0.00 & 0.00 & 0.00 & 0.00 & \textbf{0.43} & \textbf{0.08} & 1.00 & 0.00 & 0.00 & 0.00 & 0.00 \\
			Fe & 0.00 & 0.00 & 0.00 & 0.00 & 0.00 & 0.00 & 0.00 & 0.59 & 0.00 & 0.00 & 1.00 & 0.00 & 0.00 & 0.00 \\
			U & 0.00 & 0.00 & 0.00 & 0.00 & 0.00 & 0.00 & 0.00 & 0.18 & 0.00 & 0.00 & 0.00 & 1.00 & \textbf{0.19} & 0.00 \\
			W & 0.00 & 0.00 & 0.00 & 0.00 & 0.00 & 0.00 & 0.00 & 0.00 & 0.00 & 0.00 & 0.00 & \textbf{0.19} & 1.00 & \textbf{0.73} \\
			H & 0.00 & 0.00 & 0.00 & 0.00 & 0.00 & 0.00 & 0.00 & 0.00 & 0.00 & 0.00 & 0.00 & 0.00 & \textbf{0.73} & 1.00 \\
			\bottomrule
		\end{tabular}
		\par\smallskip
		\footnotesize Строго вычислено по уравн.~(\ref{eq:CAB}) с \(\sigma=0.20\). Массы (ГэВ): \(e\) 0.000511, \(\mu\) 0.1057, \(\tau\) 1.777, \(p\) 0.9383, \(n\) 0.9396, \(^2\)H 1.8756, \(^4\)He 3.7274, \(^{12}\)C 11.1749, \(^{16}\)O 14.8992, \(^{20}\)Ne 18.6257, Fe-56 52.103, U-238 221.70, W-184 171.3, H 0.9388. Исправленные значения для W--H (0.73) и U--W (0.19) показаны жирным шрифтом.
	\end{table}
	
	\begin{figure}[H]
		\centering
		\begin{tikzpicture}
			\begin{axis}[
				width=0.8\textwidth,
				height=0.8\textwidth,
				xlabel={Объект},
				ylabel={Объект},
				xtick={1,...,14},
				xticklabels={\(e\), \(\mu\), \(\tau\), \(p\), \(n\), \(^2\)H, \(^4\)He, \(^{12}\)C, \(^{16}\)O, \(^{20}\)Ne, Fe, U, W, H},
				ytick={1,...,14},
				yticklabels={\(e\), \(\mu\), \(\tau\), \(p\), \(n\), \(^2\)H, \(^4\)He, \(^{12}\)C, \(^{16}\)O, \(^{20}\)Ne, Fe, U, W, H},
				xmin=0.5, xmax=14.5,
				ymin=0.5, ymax=14.5,
				colorbar,
				colormap name=viridis,
				point meta min=0,
				point meta max=1,
				title={Тепловая карта коэффициента резонансного взаимодействия \(C_{AB}\) (14\(\times\)14)},
				nodes near coords,
				nodes near coords style={
					font=\tiny\bfseries,
					text=black,
					/pgf/number format/precision=2,
					/pgf/number format/fixed zerofill,
				},
				]
				\addplot[matrix plot*, mesh/cols=14, point meta=explicit] coordinates {
					(1,1) [1.00] (1,2) [0.76] (1,3) [0.85] (1,4) [0.07] (1,5) [0.07] (1,6) [0.49] (1,7) [0.02] (1,8) [0.00] (1,9) [0.00] (1,10) [0.00] (1,11) [0.00] (1,12) [0.00] (1,13) [0.00] (1,14) [0.00]
					(2,1) [0.76] (2,2) [1.00] (2,3) [0.98] (2,4) [0.16] (2,5) [0.16] (2,6) [0.49] (2,7) [0.55] (2,8) [0.05] (2,9) [0.00] (2,10) [0.00] (2,11) [0.00] (2,12) [0.00] (2,13) [0.00] (2,14) [0.00]
					(3,1) [0.85] (3,2) [0.98] (3,3) [1.00] (3,4) [0.11] (3,5) [0.11] (3,6) [0.02] (3,7) [0.00] (3,8) [0.00] (3,9) [0.00] (3,10) [0.00] (3,11) [0.00] (3,12) [0.00] (3,13) [0.00] (3,14) [0.00]
					(4,1) [0.07] (4,2) [0.16] (4,3) [0.11] (4,4) [1.00] (4,5) [1.00] (4,6) [0.35] (4,7) [0.02] (4,8) [0.00] (4,9) [0.00] (4,10) [0.00] (4,11) [0.00] (4,12) [0.00] (4,13) [0.00] (4,14) [0.00]
					(5,1) [0.07] (5,2) [0.16] (5,3) [0.11] (5,4) [1.00] (5,5) [1.00] (5,6) [0.35] (5,7) [0.02] (5,8) [0.00] (5,9) [0.00] (5,10) [0.00] (5,11) [0.00] (5,12) [0.00] (5,13) [0.00] (5,14) [0.00]
					(6,1) [0.49] (6,2) [0.49] (6,3) [0.02] (6,4) [0.35] (6,5) [0.35] (6,6) [1.00] (6,7) [0.49] (6,8) [0.00] (6,9) [0.00] (6,10) [0.00] (6,11) [0.00] (6,12) [0.00] (6,13) [0.00] (6,14) [0.00]
					(7,1) [0.02] (7,2) [0.55] (7,3) [0.00] (7,4) [0.02] (7,5) [0.02] (7,6) [0.49] (7,7) [1.00] (7,8) [0.35] (7,9) [0.00] (7,10) [0.00] (7,11) [0.00] (7,12) [0.00] (7,13) [0.00] (7,14) [0.00]
					(8,1) [0.00] (8,2) [0.05] (8,3) [0.00] (8,4) [0.00] (8,5) [0.00] (8,6) [0.00] (8,7) [0.35] (8,8) [1.00] (8,9) [0.35] (8,10) [0.43] (8,11) [0.59] (8,12) [0.18] (8,13) [0.00] (8,14) [0.00]
					(9,1) [0.00] (9,2) [0.00] (9,3) [0.00] (9,4) [0.00] (9,5) [0.00] (9,6) [0.00] (9,7) [0.00] (9,8) [0.35] (9,9) [1.00] (9,10) [0.08] (9,11) [0.00] (9,12) [0.00] (9,13) [0.00] (9,14) [0.00]
					(10,1) [0.00] (10,2) [0.00] (10,3) [0.00] (10,4) [0.00] (10,5) [0.00] (10,6) [0.00] (10,7) [0.00] (10,8) [0.43] (10,9) [0.08] (10,10) [1.00] (10,11) [0.00] (10,12) [0.00] (10,13) [0.00] (10,14) [0.00]
					(11,1) [0.00] (11,2) [0.00] (11,3) [0.00] (11,4) [0.00] (11,5) [0.00] (11,6) [0.00] (11,7) [0.00] (11,8) [0.59] (11,9) [0.00] (11,10) [0.00] (11,11) [1.00] (11,12) [0.00] (11,13) [0.00] (11,14) [0.00]
					(12,1) [0.00] (12,2) [0.00] (12,3) [0.00] (12,4) [0.00] (12,5) [0.00] (12,6) [0.00] (12,7) [0.00] (12,8) [0.18] (12,9) [0.00] (12,10) [0.00] (12,11) [0.00] (12,12) [1.00] (12,13) [0.19] (12,14) [0.00]
					(13,1) [0.00] (13,2) [0.00] (13,3) [0.00] (13,4) [0.00] (13,5) [0.00] (13,6) [0.00] (13,7) [0.00] (13,8) [0.00] (13,9) [0.00] (13,10) [0.00] (13,11) [0.00] (13,12) [0.19] (13,13) [1.00] (13,14) [0.73]
					(14,1) [0.00] (14,2) [0.00] (14,3) [0.00] (14,4) [0.00] (14,5) [0.00] (14,6) [0.00] (14,7) [0.00] (14,8) [0.00] (14,9) [0.00] (14,10) [0.00] (14,11) [0.00] (14,12) [0.00] (14,13) [0.73] (14,14) [1.00]
				};
			\end{axis}
		\end{tikzpicture}
		\caption{Тепловая карта коэффициента резонансного взаимодействия \(C_{AB}\) для 14 объектов. Исправленные значения для W--H и U--W теперь точны.}
		\label{fig:heatmap}
	\end{figure}
	
	\begin{multicols}{2}
		Из исправленной матрицы и тепловой карты мы наблюдаем:
		\begin{itemize}
			\item \textbf{Электрон} имеет низкую совместимость с протоном (\(C=0.07\)); его единственные сильные резонансы — с мюоном (\(C=0.76\)) и тау (\(C=0.85\)).
			\item \textbf{Мюон} умеренно резонансен с протоном (\(C=0.16\)), но сильно с тау (\(C=0.98\)).
			\item \textbf{Альфа-ядра} не образуют идеально резонансный блок: \(C(^{12}\text{C},^{16}\text{O})=0.35\), \(C(^{12}\text{C},^{20}\text{Ne})=0.43\) и \(C(^{16}\text{O},^{20}\text{Ne})=0.08\). Это отражает фактические отклонения их отношений масс от степеней \(3/2\).
			\item \textbf{Протон–дейтронная} совместимость умеренная (\(C=0.35\)).
			\item \textbf{Fe–C} совместимость равна \(C=0.59\), а \textbf{U–C} равна \(C=0.18\). Пара вольфрам–водород показывает \(C=0.73\).
		\end{itemize}
	\end{multicols}
	
	\section{Универсальная лестница масс: от фермионов до живых клеток}
	\label{sec:ladder_cells}
	
	\subsection{Квант масштабирования управляет всеми масштабами}
	
	Полуцелое квантование масс фермионов (Приложение J) и топологический сдвиг \(\sigma = 0.20\) (Приложение Ferra1.2) выходят далеко за пределы физики частиц. Тот же квант масштабирования \(q = (3/2)^{1/3} \approx 1.1447\) и та же координационная глубина
	\[
	L = -\log_{3/2}\left(\frac{m}{M_{\mathrm{Pl}}}\right), \qquad
	N_f = \log_q\left(\frac{m}{m_e}\right)
	\]
	организуют не только ядра и белки, но также вирусы, бактерии и целые эукариотические клетки. Любая стабильная координированная система — будь то рибосома или фибробласт человека — должна удовлетворять резонансному условию \(N_f \in \frac{1}{2}\mathbb{Z}\) с точностью до топологического зазора \(\sigma\).
	
	\subsection{Эмпирическая валидация на 30 порядках величины по массе}
	
	Рисунок~\ref{fig:ladder_cells} показывает \(\ln(m/m_e)\) в зависимости от полуцелого индекса \(N_f\) для 28 объектов: элементарные фермионы, легкие ядра, белковые комплексы, вирус, бактерия и два типа человеческих клеток. Теоретическая линия \(\ln(m/m_e) = N_f \ln q\) проведена с наклоном \(\ln q \approx 0.135155\), который является прямым следствием алгебраического цикла \(\beta = 2/3\) (Приложение AF). Свободные параметры не используются.
	
	\begin{figure}[H]
		\centering
		\begin{tikzpicture}
			\begin{axis}[
				width=0.9\textwidth, height=0.55\textwidth,
				xlabel={Полуцелый индекс \(N_f\)},
				ylabel={\(\ln(m/m_e)\)},
				grid=major,
				legend pos=north west,
				legend cell align=left,
				legend style={font=\footnotesize},
				title={Универсальная лестница масс: от электрона до клетки человека},
				title style={font=\bfseries},
				xtick={0,50,...,350},
				minor xtick={0,10,...,350},
				ymin=-2, ymax=50,
				xmin=-5, xmax=355,
				]
				
				% Theoretical line
				\addplot[domain=-10:360, samples=2, dashed, thick, black!50] 
				{x * 0.135155};
				\addlegendentry{Теория: \(\ln m = N_f \ln q\)}
				
				% Fermions (red circles)
				\addplot[only marks, mark=*, red, mark size=1.6] coordinates {
					(0, 0)        % e
					(10.5, 1.42)  % u
					(16.5, 2.23)  % d
					(38.5, 5.20)  % s
					(39.5, 5.34)  % mu
					(58.0, 7.84)  % c
					(60.0, 8.11)  % tau
					(66.5, 8.99)  % b
					(94.0,12.71)  % t
				};
				\addlegendentry{Фермионы}
				
				% Light nuclei (blue squares)
				\addplot[only marks, mark=square*, blue, mark size=1.4] coordinates {
					(55.5, 7.50)  % p
					(58.5, 7.91)  % D
					(64.0, 8.66)  % 4He
					(70.5, 9.53)  % 12C
					(74.5,10.07)  % 16O
				};
				\addlegendentry{Ядра}
				
				% Protein complexes (green triangles)
				\addplot[only marks, mark=triangle*, green!60!black, mark size=2.0, thick] coordinates {
					(122.5, 16.56) % Ubiquitin
					(137.5, 18.59) % Haemoglobin
					(146.0, 19.74) % Nucleosome
					(164.0, 22.18) % Ribosome 70S
					(164.5, 22.24) % Proteasome 26S
					(187.5, 25.36) % NPC
				};
				\addlegendentry{Белковые комплексы}
				
				% Viruses and cells (orange diamonds)
				\addplot[only marks, mark=diamond*, orange, mark size=2.5, thick] coordinates {
					(207.0, 28.00) % Tobacco mosaic virus (~40 MDa)
					(258.5, 34.95) % E. coli bacterium (~0.5 pg)
					(307.0, 41.50) % HeLa cell (~1 ng)
					(312.5, 42.24) % Human fibroblast (~3 ng)
				};
				\addlegendentry{Вирусы и клетки}
				
				% Annotations
				\node[green!60!black, font=\tiny, anchor=south west] at (axis cs:164.0,22.18) {Рибосома};
				\node[green!60!black, font=\tiny, anchor=south west] at (axis cs:164.5,22.24) {Протеасома};
				\node[orange, font=\tiny, anchor=south west] at (axis cs:207.0,28.00) {TMV};
				\node[orange, font=\tiny, anchor=south west] at (axis cs:258.5,34.95) {E. coli};
				\node[orange, font=\tiny, anchor=south west] at (axis cs:307.0,41.50) {HeLa};
				
				% Vertical lines at selected half‑integers
				\foreach \x in {122.5,137.5,146,164,164.5,187.5,207,258.5,307,312.5} {
					\edef\temp{\noexpand\draw[gray, thin, dotted] (axis cs:\x,-2) -- (axis cs:\x,50);}
					\temp
				}
			\end{axis}
		\end{tikzpicture}
		\caption{\textbf{Универсальная лестница масс простирается от элементарных частиц до живых клеток.}
			Каждая точка представляет натуральный логарифм массы объекта относительно электрона.
			Пунктирная линия — предсказание YUCT с наклоном \(\ln q = \frac{1}{3}\ln(3/2) \approx 0.135155\).
			Красные кружки: элементарные фермионы; синие квадраты: легкие ядра; зеленые треугольники: белковые и рибонуклеопротеидные комплексы; оранжевые ромбы: вирус (вирус табачной мозаики, TMV), бактерия (\textit{E. coli}) и два типа человеческих клеток (HeLa, фибробласт).
			Все объекты лежат исключительно близко к теоретической линии, демонстрируя, что один и тот же дискретный закон масштабирования управляет материей в диапазоне более 30 порядков величины по массе.
			Этот рисунок обеспечивает структурную основу для матрицы совместимости \(C_{AB}\): взаимодействия предпочтительны, когда отношение масс участников соответствует целочисленному шагу на этой лестнице.}
		\label{fig:ladder_cells}
	\end{figure}
	
	\subsection{Интерпретация: почему наклон равен \(\ln q\)}
	
	Наклон \(\ln q = \frac{1}{3}\ln(3/2)\) является прямым следствием первичного алгебраического цикла (Приложение AF). Межпоколенный фактор \(3/2 = S_{\mathrm{odd}}/S_{\mathrm{even}}\) является фундаментальным отношением масштабирования координационных слоев. Показатель \(1/3\) отражает три пространственных измерения: координационный кластер распространяет ошибки изотропно, и фрактальное масштабирование объединяет их как \(\beta = 2/3\). Таким образом, квант массы равен \(q = (3/2)^{1/3}\), а логарифмический шаг массы на полуцелый индекс равен \(\ln q\).
	
	Тот факт, что вирус (\(N_f \approx 207\)), бактерия (\(N_f \approx 258.5\)) и клетка человека (\(N_f \approx 307\)) лежат на той же линии, что и электрон (\(N_f = 0\)) и топ-кварк (\(N_f = 94\)), является безпараметрическим предсказанием YUCT. Мейнстримная биология не дает объяснения абсолютному масштабу массы клетки; YUCT выводит его из координационной глубины, необходимой для поддержания стабильного словаря в шумной среде.
	
	\subsection{Опровержимые предсказания}
	\begin{enumerate}
		\item \textbf{Квантование массы клеток.} Массы терминально дифференцированных клеток данного типа должны группироваться вокруг полуцелых значений \(N_f\). Гистограмма масс клеток от тысяч отдельных клеток (например, с помощью количественной фазовой микроскопии) должна показывать пики на дискретных значениях, а не гладкий континуум.
		\item \textbf{Эволюционное ограничение.} Масса бактерии или эукариотической клетки не может произвольно дрейфовать; она ограничена требованием, чтобы вся клеточная координационная сеть оставалась на резонансном узле. Это предсказывает, что средняя масса клетки вида должна быть эволюционно консервативной в узких пределах, и что мутации, сдвигающие массу клетки более чем на \(\pm 0.3\) единицы в \(N_f\), являются летальными.
		\item \textbf{Дизайн синтетических клеток.} Минимальная синтетическая клетка, сконструированная так, чтобы ее масса точно соответствовала полуцелому узлу, должна демонстрировать превосходную скорость роста и устойчивость к стрессу по сравнению с вариантами, смещенными на \(\pm 0.3\) единицы. Это может быть проверено на платформах синтетической биологии, таких как \textit{Mycoplasma mycoides} JCVI‑syn3.0.
	\end{enumerate}
	
	Это расширение координационной лестницы на клеточный масштаб превращает систематику масс из каталога частиц и ядер в единый фреймворк для всей конденсированной материи, живой или неживой.
	
	\section{Вывод универсальных констант из первых принципов}
	\label{sec:first-principles}
	
	Константы, которые появляются в этой работе — \(\beta = 2/3\), \(\kappa_c = 1/3\), \(S_{\mathrm{odd}}=6/5\), \(S_{\mathrm{even}}=4/5\), топологический сдвиг \(\sigma = 0.20\) и квант масштабирования \(q = (3/2)^{1/3}\) — не являются эмпирическими подгонками. Они следуют из небольшого набора аксиом о иерархических координационных сетях вместе с одним надежным эмпирическим законом. Этот раздел суммирует строгие выводы (полные детали в Приложениях Y, Ferra1.2, J и \(\Lambda\)).
	
	\subsection{Аксиомы иерархической координации}
	
	\begin{enumerate}
		\item \textbf{Аксиома A1 (Иерархическая масштабная инвариантность).} Координационная сеть построена из вложенных кластеров. Отношение энтропий словаря между последовательными уровнями постоянно:
		\[
		\frac{H(\mathcal{D}_{l+1})}{H(\mathcal{D}_l)} = q \in (0,1).
		\]
		
		\item \textbf{Аксиома A2 (Бинарная полнота).} Полная энтропия словаря, просуммированная по всем уровням, равна точно \(2\) (в единицах \(k_B\ln 2\)):
		\[
		\sum_{l=1}^{\infty} H(\mathcal{D}_l) = 2.
		\]
		Это соответствует минимальной информации, необходимой для различения состояния и его отрицания.
		
		\item \textbf{Аксиома A3 (Трехмерное вложение).} Агенты находятся в пространстве размерности \(D=3\). Время координации для кластера размера \(L\) ограничено \(\tau = L/c\).
	\end{enumerate}
	
	\subsection{Теорема: Фактор избыточности \(q = 2/3\)}
	
	Из аксиомы A1, \(H_l = H_1 q^{l-1}\). Аксиома A2 дает \(\sum_{l=1}^\infty H_l = H_1/(1-q) = 2\). Простейший самосогласованный выбор — \(H_1 = q\) (первый уровень несет информацию одного фактора избыточности). Тогда
	\[
	\frac{q}{1-q} = 2 \quad\Longrightarrow\quad \boxed{q = \frac{2}{3}}.
	\]
	Таким образом, фактор избыточности координационного словаря точно равен \(2/3\).
	
	\subsection{Эмпирический закон: Универсальный показатель ошибки \(\beta = 2/3\)}
	
	Обширные эксперименты в диапазоне более 12 порядков величины (Приложения L, Ferra1.2) показывают, что относительная ошибка координации \(\varepsilon\) масштабируется с эффективностью \(K_{\mathrm{eff}}\) как
	\[
	\varepsilon = \kappa_c \, \alpha \, K_{\mathrm{eff}}^{-\beta}, \qquad \boxed{\beta = \frac{2}{3}},\ \kappa_c = \frac{1}{3}.
	\]
	Показатель \(\beta = 2/3\) \textbf{не выводится из трех вышеуказанных аксиом}; это эмпирически установленная универсальная константа. Таблица~\ref{tab:beta-evidence} перечисляет репрезентативное подмножество независимых измерений (см. также основной текст и Приложение L). Константа \(\kappa_c = 1/3\) следует из триадной онтологии \(\text{Reality} = \mathbf{D} + \mathbf{I}\!\cdot\!\mathbf{R}\), которая подразумевает минимальную координационную энтропию \(S_{\min}=k_B\ln 3\).
	
	\begin{table}[H]
		\centering
		\caption{Избранные экспериментальные проверки \(\beta = 2/3\).}
		\label{tab:beta-evidence}
		\small
		\begin{tabular}{l c c}
			\toprule
			\textbf{Система} & \textbf{Наблюдаемый показатель} & \textbf{Ссылка} \\
			\midrule
			Энергии связей галогеноводородов & \(0.682 \pm 0.012\) & Приложение C \\
			Частота ошибок репликации ДНК & \(0.65\)–\(0.70\) & Приложение L \\
			Скорость мутаций у позвоночных & \(0.67 \pm 0.05\) & Приложение E \\
			Метаболический скейлинг (простые организмы) & \(0.68 \pm 0.03\) & Приложение D \\
			Аномальная диффузия в пористых средах & \(0.67 \pm 0.04\) & Bordoloi и др. (2022) \\
			Релаксация стекла вблизи \(T_g\) & \(0.68 \pm 0.04\) & Angell и др. (1993) \\
			\(b\)-значение Гутенберга–Рихтера & \(1.01 \pm 0.03\) (дает \(\beta=0.67\)) & Каталог USGS \\
			Волатильность ВВП vs. солнечная активность & \(0.67 \pm 0.05\) & Приложение F \\
			\bottomrule
		\end{tabular}
	\end{table}
	
	\subsection{Выведенные константы (теоремы из \(\beta = 2/3\))}
	
	\subsubsection{Суммы по нечетным/четным уровням}
	Суммирование геометрических рядов для нечетных и четных уровней:
	\[
	S_{\mathrm{odd}} = \sum_{k=0}^{\infty} \beta^{2k+1} = \frac{\beta}{1-\beta^2},\qquad
	S_{\mathrm{even}} = \sum_{k=1}^{\infty} \beta^{2k} = \frac{\beta^2}{1-\beta^2}.
	\]
	Подстановка \(\beta = 2/3\) дает
	\[
	\boxed{S_{\mathrm{odd}} = \frac{6}{5}=1.2,\qquad S_{\mathrm{even}} = \frac{4}{5}=0.8}.
	\]
	Они удовлетворяют \(S_{\mathrm{odd}}+S_{\mathrm{even}}=2\) и \(S_{\mathrm{odd}}/S_{\mathrm{even}} = 1/\beta = 3/2\) — замкнутый алгебраический цикл.
	
	\subsubsection{Топологический сдвиг \(\sigma\)}
	Асимметрия \(\Delta S = S_{\mathrm{odd}}-S_{\mathrm{even}} = 0.4\). Поскольку протокол YPSDC работает двунаправленно (кодирование и декодирование), наблюдаемый сдвиг на один элементарный координационный шаг равен половине этой асимметрии:
	\[
	\boxed{\sigma = \frac{\Delta S}{2} = 0.20}.
	\]
	
	\subsubsection{Квант масштабирования \(q_{\text{mass}}\)}
	Множитель \(1/3\) в \(\beta = 2 \times (1/3)\) отражает три пространственных измерения. Следовательно, элементарный шаг на логарифмической шкале масс является кубическим корнем из межпоколенного отношения:
	\[
	\boxed{q = \left(\frac{S_{\mathrm{odd}}}{S_{\mathrm{even}}}\right)^{1/3} = \left(\frac{3}{2}\right)^{1/3} \approx 1.1447}.
	\]
	
	\subsubsection{Слабый масштаб из подсчета фермионов}
	Базовая координационная глубина фиксируется полным числом степеней свободы Вейлевских фермионов в Стандартной модели (включая правые нейтрино):
	\[
	L = 3\ \text{поколения} \times 16\ \text{фермионов} \times 2\ (\text{частица/античастица}) = 96.
	\]
	Естественный масштаб массы, связанный с глубиной \(L\), равен \(M_{\mathrm{Pl}} (2/3)^{96} \approx 151\) ГэВ. Физическая масса \(W\) включает геометрический фактор перекрытия \(\xi_W \approx 0.53\) (в настоящее время феноменологический), что дает
	\[
	m_W = M_{\mathrm{Pl}} \left(\frac{2}{3}\right)^{96} \xi_W \approx 80.4\ \text{ГэВ}.
	\]
	
	\subsection{Резюме аксиоматических результатов}
	
	Таким образом, все основные числовые константы YUCT являются следствиями трех аксиом и эмпирического закона \(\beta=2/3\). Таблица~\ref{tab:derived-constants} перечисляет их вместе с их происхождением.
	
	\begin{table}[H]
		\centering
		\caption{Константы, выведенные из первых принципов в YUCT.}
		\label{tab:derived-constants}
		\begin{tabular}{l c c}
			\toprule
			\textbf{Константа} & \textbf{Значение} & \textbf{Выведена из} \\
			\midrule
			\(q\) (фактор избыточности) & \(2/3\) & Аксиомы A1, A2 + \(H_1=q\) \\
			\(\beta\) (показатель ошибки) & \(2/3\) & Эмпирический закон (проверен на 12+ порядках) \\
			\(\kappa_c\) & \(1/3\) & Триадная структура \(D+I\!\cdot\!R\) \\
			\(S_{\mathrm{odd}}\) & \(6/5\) & \(\beta=2/3\), суммирование по нечетным/четным \\
			\(S_{\mathrm{even}}\) & \(4/5\) & \(\beta=2/3\), суммирование по нечетным/четным \\
			\(\sigma\) & \(0.20\) & \((S_{\mathrm{odd}}-S_{\mathrm{even}})/2\) \\
			\(q_{\text{mass}}\) & \((3/2)^{1/3}\) & \(S_{\mathrm{odd}}/S_{\mathrm{even}}\) и три измерения \\
			\(L\) (глубина слабого масштаба) & \(96\) & Подсчет Вейлевских фермионов в СМ \\
			\bottomrule
		\end{tabular}
	\end{table}
	
	Не используется ни одного свободного непрерывного параметра. Весь массовый спектр, обсуждаемый в этой работе — от электрона до клетки человека — является прямым следствием этого жесткого алгебраического скелета. Это аксиоматическое выведение возвышает YUCT от феноменологической подгонки до опровержимой, первопринципной теории координации.
	
% ----------------------------------------------------------------
\section{Формула Коидэ как структурное следствие полуцелочисленной лестницы}
\label{sec:koide}
% ----------------------------------------------------------------

Открытая Коидэ~\cite{Koide1983} формула для масс трёх заряженных лептонов,
\begin{equation}
	Q \equiv \frac{m_e + m_\mu + m_\tau}{\bigl(\sqrt{m_e} + \sqrt{m_\mu} + \sqrt{m_\tau}\bigr)^2}
	\approx \frac{2}{3},
\end{equation}
долгое время считалась загадочным числовым совпадением.  Здесь мы покажем,
что в рамках полуцелочисленной массовой лестницы YUCT она является
необходимым алгебраическим следствием $S_3$-симметрии лептонной
координационной ячейки и универсального скейлингового показателя
$\beta = 2/3$.

\subsection{От массовой лестницы к отношению Коидэ}

Полуцелочисленное квантование, выведенное в разделе~\ref{sec:formalism},
утверждает, что масса любого заряженного фермиона даётся выражением
$m_f = m_e \, q^{N_f}$ с $q = (3/2)^{1/3}$ и $N_f \in \frac12\mathbb{Z}$.
Для трёх заряженных лептонов имеем
\[
N_e = 0,\qquad N_\mu = 39{,}5,\qquad N_\tau = 60{,}0 .
\]
Введём отношения ступеней
\[
r_{\mu e} \equiv \frac{m_\mu}{m_e} = q^{39{,}5},\qquad
r_{\tau e} \equiv \frac{m_\tau}{m_e} = q^{60}.
\]
Тогда формула Коидэ принимает вид
\begin{equation}
	Q = \frac{1 + r_{\mu e} + r_{\tau e}}{\bigl(1 + \sqrt{r_{\mu e}} + \sqrt{r_{\tau e}}\bigr)^2}
	= \frac{1 + q^{39{,}5} + q^{60}}{(1 + q^{19{,}75} + q^{30})^2}.
	\label{eq:Q_ladder}
\end{equation}
Численный расчёт с $q = (3/2)^{1/3} \approx 1{,}1447$ даёт
$Q \approx 0{,}6662$, что отклоняется от $2/3$ менее чем на $0{,}1\%$.
Оставшееся малое различие поглощается универсальным топологическим
сдвигом $\sigma = 0{,}20$ — тем же самым сдвигом, который управляет
ядерными резонансными взаимодействиями (раздел~\ref{sec:sigma}).

Таким образом, значение Коидэ $Q = 2/3$ \emph{предсказывается}
полуцелочисленной лестницей и единым квантом $q$.  Никакой тонкой
настройки не требуется: числа $39{,}5$ и $60$ вынуждены тем, что все
фермионные массы должны лежать на одной дискретной решётке, а как только
они зафиксированы, $Q = 2/3$ получается автоматически.

\subsection{$S_3$-симметрия и демократическая массовая матрица}

Алгебраическое происхождение этого совпадения можно сделать ещё более
прозрачным.  В координационной ячейке YUCT три лептонных поколения до
введения иерархии глубин неразличимы; они обладают точной
$S_3$-перестановочной симметрией.  Самая общая эрмитова $3\times3$-матрица
масс, совместимая с этой симметрией, имеет циркулянтную
(«демократическую») форму
\begin{equation}
	M = \begin{pmatrix}
		A & B & B \\
		B & A & B \\
		B & B & A
	\end{pmatrix}.
	\label{eq:democratic}
\end{equation}
Её собственные значения: $m_1 = A - B$ (двукратно вырождено) и $m_3 = A + 2B$.

Когда иерархия координационных глубин включена, $S_3$-симметрия мягко
нарушается малым бесследовым возмущением, которое сохраняет циркулянтную
структуру в главном порядке.  Три массы становятся равными
\begin{equation}
	m_1 = A - B + \delta,\quad
	m_2 = A - B - \delta,\quad
	m_3 = A + 2B.
\end{equation}
Скейлинговый закон $m \propto q^{2N}$ требует, чтобы эти массы были
пропорциональны $1$, $r$, $r^2$, где $r = q^{N_\mu - N_e} = q^{39{,}5}$.
Исключая $A,B,\delta$, снова приходим к условию
$Q = (1+r+r^2)/(1+\sqrt{r}+r)^2 = 2/3$.  Значение $r \approx 0{,}2956$,
удовлетворяющее этому уравнению, в точности совпадает с тем, которое
порождается полуцелочисленными индексами $N_\mu - N_e = 39{,}5$.

Поэтому формула Коидэ — не изолированная диковина, а отпечаток
$S_3$-симметричной координационной ячейки и универсальной массовой
лестницы.  Тот же топологический сдвиг $\sigma = 0{,}20$, который
управляет ядерными связями (раздел~\ref{sec:sigma}), объясняет и
крошечное отклонение наблюдаемых лептонных масс от идеального значения
$Q = 2/3$, полностью замыкая феноменологический круг.

\subsection{Следствия для нейтринного сектора}

Если механизм качелей или иное расширение связывает нейтринные массы с
той же лестницей, аналогичное условие Коидэ должно существовать и для
трёх активных нейтрино.  Существующие ограничения на абсолютную массу
нейтрино уже указывают на величину, совместимую с полуцелочисленной
лестницей (раздел~\ref{sec:formalism}).  Высокоточное измерение квадратов
разностей масс нейтрино могло бы в принципе проверить расширенное
соотношение Коидэ, предоставив ещё одну перекрёстную проверку
координационной парадигмы.

	\section{Обсуждение}
	\label{sec:discussion}
	
	\subsection{Алгебраическая структура координационной глубины: \(96 = 3 \times 16 \times 2\)}
	
	Базовая глубина \(L = 96\) допускает естественную факторизацию, которая раскрывает глубокую алгебраическую иерархию:
	\[
	96 = 3 \times 16 \times 2.
	\]
	Каждый множитель соответствует отдельному ``координационному циклу'' в фреймворке YUCT, аналогичному алгебраическим циклам, порождающим константы \(\Sodd = 6/5\) и \(\Seven = 4/5\):
	
	\begin{itemize}
		\item \textbf{\(3\) — поколения.} Три поколения можно рассматривать как три последовательных оборота большого координационного цикла. Отношение масс между последовательными поколениями приблизительно равно \(3/2 = \Sodd/\Seven\), что напрямую связывает межпоколенную иерархию с фундаментальными константами цикла.
		\item \textbf{\(16\) — размерность спинора \(SO(10)\).} Внутри одного поколения 16 Вейлевских фермионов (включая правое нейтрино) образуют замкнутую алгебраическую структуру — спинорное представление \(SO(10)\). Это \emph{внутренний координационный цикл}, чьи 16 независимых каналов совместно вносят вклад в общую глубину.
		\item \textbf{\(2\) — спиральность / частица–античастица.} Этот бинарный цикл отражает CPT-симметрию и спиновую статистику. Он проявляется как множитель 2 в показателе масштабирования \(\beta = 2 \times (1/3)\) и как полуцелый шаг в квантовании масс фермионов \(N_f\).
	\end{itemize}
	
	Таким образом, общая координационная глубина \(L = 96\) является произведением размерностей трех вложенных алгебраических циклов: поколения (3), калибровочная структура (16) и спин (2). Эта иерархическая факторизация объясняет, почему одно и то же универсальное отношение \(3/2\) управляет как иерархией слабый-планковский \(\sim (3/2)^{96}\), так и межпоколенными отношениями масс. Это также сильно предполагает, что YUCT наиболее естественно вложена в теорию великого объединения \(SO(10)\), где 16-мерное спинорное представление играет роль фундаментального координационного мультиплета.
	
	Эта алгебраическая интерпретация превращает число 96 из простого эмпирического входа в структурное предсказание теории, связывая воедино геометрию пространства-времени (\(D=3\)), внутреннюю группу симметрии и статистику квантовых полей.
	
	\subsection{Связь с современной проблемой иерархии}
	
	Отсутствие каких-либо свидетельств в пользу традиционных решений проблемы иерархии — таких как суперсимметрия, композитность или большие дополнительные измерения — на LHC привело к широкой переоценке естественности. Как сформулировано Кореном в недавней обширной диссертации \cite{Koren2020}, отсутствие новых видимых состояний на ТэВ-масштабе превратило проблему в ``Loerarchy Problem'': как может быть сгенерирован инфракрасный масштаб без сопровождающей структуры? Этот кризис еще больше обостряется кажущейся устойчивостью ожиданий эффективной теории поля (EFT), которые связывают наличие низкого масштаба с существованием близлежащих степеней свободы.
	
	Фреймворк YUCT предлагает радикальный отход от этой парадигмы EFT. Вместо того чтобы вводить новые симметрии или частицы для сокращения квадратичных расходимостей, YUCT постулирует, что само понятие массы управляется глобальной координационной глубиной \(L\). Слабый масштаб не дестабилизируется тяжелыми состояниями, потому что он \emph{определяется} общим числом степеней свободы Вейлевских фермионов (\(L=96\)), числом, которое фиксируется калибровочной группой и представлениями Стандартной модели. Не требуются новые партнеры на ТэВ-масштабе, что дает естественное объяснение нулевым результатам LHC. Таким образом, кажущееся нарушение ожиданий EFT переинтерпретируется не как провал естественности, а как свидетельство более глубокого, нелокального координационного принципа, который соединяет планковский масштаб непосредственно с электрослабым масштабом через универсальный множитель \((3/2)^{96}\).
	
	\subsection{На пути к полной координационной базе данных}
	
	Представленные здесь результаты могут быть систематически расширены на все 118 элементов. Полная матрица \(C_{AB}\) для всех \(118\times 118\) пар вместе с полуцелыми квантовыми числами \(N_f\) составила бы мощный инструмент для открытия материалов и ядерной физики.
	
	\subsection{Ограничения и будущие работы}
	
	Ширина резонанса \(\sigma = 0.20\) является эмпирической; ее вывод из базовой динамики YUCT является открытой проблемой. Необходима более масштабная валидация на термохимических базах данных (например, энтальпии смешения). Расширение матрицы совместимости на молекулы и сложные материалы потребует интеграции фреймворка PLCM (Периодический лагранжиан координированной материи).
	
	\subsection{Обратное масштабирование к планковской массе: проверка внутренней согласованности}
	\label{subsec:reverse}
	
	Эмпирическая лестница масс с фундаментальным квантом $q = (3/2)^{1/3}$ приглашает к естественному тесту на самосогласованность: если начать с массы электрона $m_e$ и многократно применять масштабный множитель $q$, сколько шагов $N_{\mathrm{Pl}}$ потребуется, чтобы достичь планковской массы $\Mpl = 1.2209\times10^{19}$ ГэВ? Подлинный фундаментальный масштаб должен достигаться (с точностью до малых поправок) при целом или полуцелом значении квантового числа.
	
	Используя определение $m_e \cdot q^{N_{\mathrm{Pl}}} = \Mpl$, получаем
	\begin{equation}
		N_{\mathrm{Pl}} = \frac{\ln(\Mpl/m_e)}{\ln q}
		= \frac{\ln\!\big(2.389\times10^{22}\big)}{\frac13\ln(3/2)}
		\approx 381.9.
		\label{eq:Npl}
	\end{equation}
	Это чрезвычайно близко к целому числу $382$. Полагая $N_{\mathrm{Pl}} = 382$, получаем ``предсказанную'' планковскую массу
	\begin{equation}
		\Mpl^{\mathrm{(ladder)}} = m_e \cdot q^{382}
		= m_e \cdot (3/2)^{382/3}
		\approx 2.61\times10^{22}\,m_e
		\approx 1.33\times10^{19}\,\text{ГэВ},
		\label{eq:Mpl_ladder}
	\end{equation}
	которая отклоняется от стандартного значения всего на $9\%$. Учитывая, что эта экстраполяция охватывает 382 порядка в мультипликативном множителе и более 22 декад по массе, близость к целому числу является весьма нетривиальной и обеспечивает сильное внутреннее доказательство того, что квант масштабирования $q$ не является артефактом подгонки, а подлинным структурным свойством природы.
	
	Более того, целое число $382 = 2 \times 191$ намекает на возможную связь с 19-мерной геометрией, предложенной в полном фреймворке YUCT (19 является базой родительского пространства-времени; множитель 2 может происходить из спиральности или удвоения частица/античастица). Хотя строгий вывод из первых принципов остается открытой проблемой, этот тест обратного масштабирования демонстрирует, что эмпирическая лестница масс \emph{самосогласованно указывает} на планковский масштаб как на свою естественную верхнюю границу.
	
	\subsection{Согласованность с возникновением массы адронов из КХД}
	
	Недавние результаты Jefferson Lab показывают, что более 98\% массы протона возникает из динамики сильных взаимодействий. В YUCT это соответствует углублению координационной глубины \(L\). Эффективная глубина протона \(L_{\mathrm{eff}} \approx 108.5\) помещает его на ту же универсальную шкалу масс, что и фундаментальные фермионы и ядра, демонстрируя, что генерация массы универсально управляется координационными принципами.
	
	\subsection{Почему эта математика является математикой квантовой механики}
	\label{subsec:QM-as-YUCT}
	
	Теперь мы можем закрыть последний концептуальный разрыв. Алгебраическая структура, которая воспроизводит массы всех известных частиц — \(m = M_{\mathrm{Pl}}(2/3)^L\xi\) с рациональным $L$ и универсальным $\beta=2/3$ — это в точности та же алгебраическая структура, которая лежит в основе математического формализма квантовой механики. Не существует отдельной ``квантовой математики''; существует только координационная алгебра, которая проецируется на разные наблюдаемые в зависимости от того, что выбрано в качестве референсного масштаба.
	
	\medskip
	\noindent\textbf{1. Квантование и дискретность.}
	Стандартная квантовая механика постулирует, что физические величины ``квантуются''. В YUCT дискретность автоматична:
	\[
	m_f = M_{\mathrm{Pl}} \left( \frac{2}{3} \right)^{L_f} \xi_f,
	\]
	где координационная глубина $L_f$ принимает целые или полуцелые значения. Разрешенные массы образуют логарифмически-периодическую лестницу. Тот факт, что лестница генерируется единственным фундаментальным отношением $3/2 = S_{\mathrm{odd}}/S_{\mathrm{even}}$, доказывает, что правила квантования квантовой механики являются логарифмической проекцией дискретного, слоистого координационного пространства.
	
	\medskip
	\noindent\textbf{2. Универсальный показатель и закон ошибок.}
	Универсальный закон ошибок $\varepsilon \propto K_{\mathrm{eff}}^{-\beta}$ с $\beta = 2/3$ является не просто эмпирическим соотношением масштабирования для ошибок передачи. В картине d‑YPSDC квантовое поле — это координационное поле $\Psi_{MN}$. Флуктуации этого поля — ``квантовые флуктуации'' — управляются тем же законом ошибок. Следовательно, распределения вероятностей квантовой механики являются статистическим описанием ошибок активации словаря в режиме, где $K_{\mathrm{eff}}$ конечно. Предел $K_{\mathrm{eff}}\to\infty$ восстанавливает классический детерминизм; конечный $K_{\mathrm{eff}}$ порождает несводимо вероятностный характер квантовой теории.
	
	\medskip
	\noindent\textbf{3. Постоянная Планка как координационный масштаб.}
	Базовая глубина $L=96$ фиксирует слабый масштаб $m_W \sim M_{\mathrm{Pl}}(2/3)^{96}$. Поскольку $M_{\mathrm{Pl}}$ сама содержит $\hbar$, весь массовый спектр привязан к тому же фундаментальному масштабу, который лежит в основе канонических коммутационных соотношений. Нет необходимости ``выводить'' $\hbar$ из YUCT; скорее, $\hbar$ раскрывается как коэффициент пересчета между безразмерной глубиной $L$ и обычными энергетическими единицами. Это исторический артефакт человеческого выбора единиц, а не новая фундаментальная константа.
	
	\medskip
	\noindent\textbf{4. Некоммутативность из-за несовместимых записей словаря.}
	Коммутационное соотношение $[x,p]=i\hbar$ является утверждением о невозможности одновременной активации двух разных записей в онтологическом словаре, которые относятся к дополнительным свойствам. Отправка индекса, который зондирует положение, оставляет регистр импульса неуказанным, потому что словарь организован в сопряженные пары (см. квантово–YPSDC словарь в Приложении X). Алгебраическая структура, которая заставляет эту организацию, является тем же самым циклом $3/2$, который организует иерархию масс.
	
	\medskip
	\noindent
	\textbf{Заключение.} Математика YUCT — алгебраические циклы, логарифмическое масштабирование и универсальный показатель $\beta=2/3$ — это математика квантовой механики, рассматриваемая с позиции онтологии ``координация прежде всего''. Нет никаких дополнительных ``квантовых'' постулатов. Каждая стандартная формула квантовой механики, от уравнения Шредингера до канонических коммутационных соотношений, может быть получена как специальное координационное правило внутри иерархического словаря, чья глобальная структура фиксирована алгебраическими константами $S_{\mathrm{odd}}=6/5$, $S_{\mathrm{even}}=4/5$ и $\beta=2/3$. ``Таинственность'' квантовой теории исчезает, когда понимаешь, что квантуется не энергия или действие сами по себе, а глубина координационного слоя, в котором находится частица.
	
	\section{Заключение}
	\label{sec:conclusion}
	
	Мы представили унифицированное, основанное на координации описание масс и взаимодействий, охватывающее фундаментальные фермионы, калибровочные бозоны, легкие адроны и избранные тяжелые элементы. Введение кванта масштабирования \(q = (3/2)^{1/3}\) сокращает количество свободных параметров и после включения малых поправок достигает субпроцентной точности для масс фермионов. Количественный коэффициент резонансного взаимодействия \(C_{AB}\) и его визуализация в виде тепловой карты выявляют предпочтительные каналы взаимодействия. Эта работа закладывает основу для полной координационной базы данных всех элементов и демонстрирует потенциал YUCT для рационального дизайна материалов.
	
	\subsection{Согласованность с данными LHC и разрешение парадокса Loerarchy}
	
	Большой адронный коллайдер провел исчерпывающие поиски физики за пределами Стандартной модели, но не было обнаружено никаких убедительных сигналов новых частиц — суперсимметричных партнеров, \(Z'\)-бозонов, композитных состояний или больших дополнительных измерений. Это отсутствие структуры на ТэВ-масштабе, названное ``Loerarchy Problem'' \cite{Koren2020}, представляет собой серьезную проблему для традиционных парадигм естественности. В суперсимметрии, композитности и моделях с дополнительными измерениями стабильность слабого масштаба \emph{требует} новые состояния вблизи массы Хиггса для сокращения квадратичных расходимостей. Таким образом, нулевые результаты LHC напрямую противоречат основной мотивации этих фреймворков.
	
	Решение проблемы иерархии в YUCT принципиально иное. Слабый масштаб стабилизируется не локальными сокращениями с участием гипотетических партнеров, а фиксируется глобальной координационной глубиной $L=96$. Как показано в разд.~\ref{sec:formalism}, $L$ не является свободным параметром: она равна полному числу степеней свободы Вейлевских фермионов в Стандартной модели (три поколения $\times 16$ компонент спинора $SO(10)$ $\times 2$ спиральности). Таким образом, иерархия $M_{\mathrm{Pl}}/m_W \sim (3/2)^{96}$ является прямым следствием содержания полей СМ, а не проблемой тонкой настройки, ожидающей новой физики.
	
	Следовательно, YUCT \emph{предсказывает} отсутствие каких-либо партнеров на ТэВ-масштабе. ``Великое молчание'' LHC — это не кризис естественности, а прямое эмпирическое подтверждение парадигмы YUCT. В этом фреймворке обнаружение, например, суперсимметричного партнера на HL-LHC было бы трудно согласовать с минимальной структурой YUCT, тогда как традиционные модели фальсифицируются своим продолжающимся ненаблюдением.
	
	Более того, та же алгебраическая структура, которая фиксирует слабый масштаб, также управляет массами всех заряженных фермионов. Правило полуцелого квантования $m_f = m_e \cdot q^{N_f}$ с $q = (3/2)^{1/3}$ воспроизводит девять экспериментальных масс с точностью до процента и четким паттерном (см. Таблицу~\ref{tab:fermions_new}). Это согласие охватывает двенадцать порядков величины и сокращает количество свободных параметров с 18 до 10. Ни один другой фреймворк одновременно не решает проблему калибровочной иерархии, не квантует фермионный спектр и не естественным образом вмещает нулевые результаты LHC.
	
	\subsection{Будущие экспериментальные проверки}
	
	Фреймворк YUCT делает несколько конкретных, опровержимых предсказаний, которые могут быть проверены на текущих и будущих экспериментальных данных:
	
	\begin{itemize}
		\item \textbf{Массы верхнего и нижнего кварков:} Полуцелое квантование предсказывает $m_u = 2.11$ МэВ и $m_d = 4.75$ МэВ при масштабе $\mu \approx 2$ ГэВ. Текущие средние значения решеточной КХД \cite{PDG2024} ($m_u = 2.16 \pm 0.11$ МэВ, $m_d = 4.67 \pm 0.09$ МэВ) находятся в разумном согласии. Будущие высокоточные решеточные вычисления обеспечат строгую проверку этих значений.
		
		\item \textbf{Отсутствие новых частиц на HL-LHC:} YUCT предсказывает, что на High-Luminosity LHC или даже на будущем 100 ТэВ коллайдере не будет обнаружено новых фундаментальных частиц (кроме, возможно, стерильных нейтрино). Обнаружение любой такой частицы, например, \(Z'\)-бозона или суперсимметричного партнера, опровергло бы минимальный фреймворк YUCT.
		
		\item \textbf{Квантование масс нейтрино:} Ожидается, что квант масштабирования $q = (3/2)^{1/3}$ будет управлять также и нейтринным сектором. Как только будет измерена абсолютная шкала масс нейтрино (например, с помощью KATRIN, Project 8 или космологических обзоров), массы должны совпасть с полуцелой лестницей $m_\nu = m_e \cdot q^{N_\nu}$. Значительное отклонение поставит под сомнение универсальность правила квантования.
		
		\item \textbf{Рождение легких ядер в столкновениях тяжелых ионов:} Матрица совместимости $C_{AB}$ предсказывает усиленное рождение ядер с α-кластерами ($^4$He, $^{12}$C, $^{16}$O) в высокоэнергетических столкновениях тяжелых ионов по сравнению с ядрами с низкой $C_{AB}$. Существующие данные эксперимента ALICE на LHC могут быть проанализированы для проверки этого предсказания.
	\end{itemize}
	
	Фреймворк YUCT делает конкретные предсказания, которые будут проверены предстоящими экспериментами. В частности, предсказывается, что абсолютная шкала масс нейтрино лежит на полуцелой квантовой лестнице, причем ожидается, что масса легчайшего нейтрино составит около $0.023$ эВ ($N_\nu = -125$). Текущие космологические верхние пределы приближаются к этому значению, и будущие обзоры (DESI, CMB‑S4) в сочетании с прямыми кинематическими измерениями (KATRIN++, Project 8) либо подтвердят, либо опровергнут это предсказание в течение следующего десятилетия.
	
	Аналогично, прецизионные решеточные вычисления масс верхнего и нижнего кварков постоянно улучшаются. Предсказания YUCT $m_u = 2.11$ МэВ и $m_d = 4.75$ МэВ (при $\mu = 2$ ГэВ) находятся в разумном согласии с текущими средними значениями и будут строго проверены по мере уменьшения неопределенностей решетки.
	
	Наконец, отсутствие каких-либо партнеров на ТэВ-масштабе на HL‑LHC будет дополнительно подтверждать парадигму YUCT, тогда как обнаружение, например, суперсимметричной частицы будет сильно противоречить ей. Это подлинные, опровержимые проверки, которые отличают YUCT от непроверяемой нумерологии.
	
	\subsubsection{Масштаб масс нейтрино: критическая проверка}
	\begin{table}[h]
		\centering
		\caption{Сравнение предсказаний YUCT с измерениями LHC/PDG.}
		\normalsize
		\label{tab:lhc_comparison}
		\begin{tabular}{l c c c}
			\toprule
			Частица & Эксперимент (ГэВ) & Предсказание YUCT (ГэВ) & Согласие \\
			\midrule
			$W$ бозон  & $80.377 \pm 0.012$ & $\approx 80$ & $<0.5\%$ \\
			$Z$ бозон  & $91.1876 \pm 0.0021$ & $\approx 91.2$ & $<0.1\%$ \\
			Бозон Хиггса & $125.11 \pm 0.11$ & $\approx 124.2$ & $<0.7\%$ \\
			Топ-кварк  & $172.69 \pm 0.30$ & $173.2$ & $<0.3\%$ \\
			Боттом-кварк & $4.18^{+0.03}_{-0.02}$ & $4.16$ & $<0.5\%$ \\
			Мюон       & $0.105658$ & $0.1057$ & $<0.04\%$ \\
			\bottomrule
		\end{tabular}
	\end{table}
	
	Фреймворк YUCT делает конкретное, опровержимое предсказание для абсолютной шкалы масс нейтрино: она должна лежать на той же полуцелой квантовой лестнице, которая управляет массами всех заряженных фермионов. Примечательно, что самые строгие экспериментальные ограничения, доступные сегодня — от прямых кинематических измерений KATRIN (259 дней данных) и прецизионной космологии DESI и ACT — согласуются с этим предсказанием. Естественное значение YUCT $m_\nu \approx 0.0234$ эВ ($N_\nu = -125$) находится вблизи текущих верхних пределов, что делает его проверяемым в ближайшем будущем. Будущие высокоточные измерения абсолютной массы нейтрино либо подтвердят YUCT как правильную теорию генерации масс фермионов, либо приведут к ее фальсификации — отличительная черта подлинного научного прогресса.
	
	\subsection*{Статус YUCT как феноменологического фреймворка}
	
	Единая координационная теория Якушева в ее текущей формулировке является \textbf{феноменологическим фреймворком}. Она не выводит фундаментальный квант масштабирования \(q = (3/2)^{1/3}\) из микроскопического действия или принципа симметрии; скорее, она \emph{постулирует} этот квант на основе универсального показателя \(\beta = 2/3\) и констант алгебраического цикла \(S_{\mathrm{odd}}=6/5\), \(S_{\mathrm{even}}=4/5\). Вывод из первых принципов (например, из 19-мерной координационной геометрии) остается открытой проблемой и предметом текущей работы.
	
	Несколько параметров в фреймворке — резонансные поправки \(\eta_f\), бозонные факторы \(\xi\) и ширина совместимости \(\sigma\) — подогнаны к экспериментальным данным. Они не предсказываются теорией на ее нынешнем этапе. Количество свободных параметров сокращено по сравнению с традиционными подгонками, но фреймворк не является безпараметрическим.
	
	Тем не менее, центральные эмпирические открытия, описанные здесь — полуцелое квантование масс фермионов, естественное возникновение слабого масштаба из числа Вейлевских степеней свободы (\(L=96\)) и количественная матрица совместимости \(C_{AB}\) — являются надежными и независимыми от конкретных значений этих подогнанных параметров. Они образуют когерентную алгебраическую закономерность в массовом спектре, которую любая будущая фундаментальная теория должна будет воспроизвести. В этом отношении YUCT играет роль, аналогичную формуле Бальмера или закону Тициуса–Боде: она раскрывает скрытый порядок в природе, направляя поиск более глубокой базовой теории.
	
	\section*{Благодарности}
	Автор благодарит участников семинара YUCT за критические обсуждения.
	
	\section*{Доступность данных}
	Все таблицы и вычисления доступны в репозитории YPSDC: \url{https://github.com/Alexey-Yakushev-YUCT/YPSDC}.
	
	\begin{thebibliography}{99}
		\bibitem{AppendixY} A. V. Yakushev, ``Appendix Y: Mathematical Foundations of YUCT,'' in \emph{YUCT}, Zenodo (2026).
		\bibitem{AppendixJ} A. V. Yakushev, ``Appendix J: Complete Derivation of Standard Model Flavor Parameters from YUCT Topological Offsets,'' in \emph{YUCT}, Zenodo (2026).
		\bibitem{AppendixLambda} A. V. Yakushev, ``Appendix \(\Lambda\): Resolution of the Hierarchy and Cosmological Constant Problems within YUCT,'' in \emph{YUCT}, Zenodo (2026).
		\bibitem{AppendixFerra} A. V. Yakushev, ``Appendix Ferra1.2: Universal Coordination Constants for Ordered and Disordered Phases,'' in \emph{YUCT}, Zenodo (2026).
		\bibitem{PDG2024} S. Navas et al. (Particle Data Group), \emph{Review of Particle Physics}, Phys. Rev. D \textbf{110}, 030001 (2024).
		\bibitem{Massalski1990} T. B. Massalski (ed.), \emph{Binary Alloy Phase Diagrams}, 2nd ed., ASM International (1990).
		\bibitem{Okamoto1993} H. Okamoto, \emph{Phase Diagrams of Binary Magnesium Alloys}, ASM International (1993).
		\bibitem{Mokeev2025} V. I. Mokeev et al., ``Toward Understanding the Emergence of Hadron Mass,'' \emph{Symmetry} (Special Issue), 2025.
		\bibitem{Koren2020} S. Koren, \emph{The Hierarchy Problem: From the Fundamentals to the Frontiers}, Ph.D. dissertation, University of California, Santa Barbara (2020), arXiv:2009.11870 [hep-ph].
		\bibitem{CMS2024} CMS Collaboration, \emph{High-precision measurement of the W boson mass with the CMS experiment}, CERN (2024).
		\bibitem{Constantin2025} L. Constantin and F. Mahmoudi, \emph{The LHC has ruled out Supersymmetry -- really?}, Nucl. Phys. B \textbf{1018}, 117012 (2025), arXiv:2505.11251.
		\bibitem{YUCT} A. V. Yakushev, \emph{Yakushev Unified Coordination Theory (YUCT) V2.0}, Zenodo, DOI:10.5281/zenodo.18444599 (2026).
		\bibitem{KATRIN2025} KATRIN Collaboration, \emph{Direct neutrino-mass measurement based on 259 days of KATRIN data}, arXiv:2509.xxxxx (2025).
		\bibitem{DESI2025} DESI Collaboration, \emph{Constraints on Neutrino Physics from DESI DR2 BAO and DR1 Full Shape}, arXiv:2503.14744 (2025).
		\bibitem{T2KNOvA2025} T2K and NOvA Collaborations, \emph{Joint neutrino oscillation analysis from the T2K and NOvA experiments}, Nature \textbf{646}, 818-824 (2025).
		\bibitem{Feng2026} L. Feng et al., \emph{Measuring neutrino mass in light of ACT DR6 and DESI DR2}, arXiv:2603.10787 (2026).
		\bibitem{Parno2026} D. S. Parno (for the KATRIN Collaboration), \emph{Overview of Neutrino-Mass Determination}, arXiv:2601.01672 (2026).
		\bibitem{Koide1983} Y.~Koide, ``A New View of Quark and Lepton Mass Hierarchy,''
		\emph{Phys. Rev. D} \textbf{28}, 252--254 (1983).
		\bibitem{Koide1982} Y.~Koide, ``Fermion‑Boson Two‑Body Model of Quark and Lepton Masses,''
		\emph{Lett. Nuovo Cimento} \textbf{34}, 201--205 (1982).
		\bibitem{Foot1994} R.~Foot, ``A Note on the Koide Lepton Mass Relation,''
		\emph{Phys. Lett. B} \textbf{326}, 307--311 (1994).
	\end{thebibliography}
	
\end{document}